\documentclass[11pt,a4paper,openany]{book}
\usepackage[T1]{fontenc}
\usepackage[utf8]{inputenc}
\usepackage[french]{babel}
\usepackage{newtxtext,newtxmath}
\usepackage{microtype}
\usepackage{geometry}
\geometry{margin=28mm,headheight=15pt}
\emergencystretch=1.5em
\usepackage{amsmath,amscd,mathrsfs,euscript,bbm}
\usepackage{eurosym}
\usepackage{comment,url,booktabs,array,enumitem,etoolbox}
\usepackage{xcolor}
\usepackage{pstricks,pst-plot,pst-fill,pst-node,pst-text,pst-3d}
\usepackage{fancyhdr}
\usepackage{titlesec}
\usepackage[hidelinks,pdfusetitle]{hyperref}
\usepackage{bookmark}

\definecolor{forest}{HTML}{234B3A}
\definecolor{terracotta}{HTML}{B45F3C}
\definecolor{ink}{HTML}{25302B}
\color{ink}
\hypersetup{colorlinks=true,linkcolor=forest,urlcolor=terracotta,citecolor=forest,
  pdftitle={Théorie de la décision et théorie des jeux},
  pdfauthor={Olivier Gossner},
  pdfsubject={Notes de cours},
  pdfkeywords={théorie de la décision, théorie des jeux, risque, équilibre de Nash}}

\pagestyle{fancy}
\fancyhf{}
\fancyhead[LE]{\small\color{forest}\thepage\quad\nouppercase{\leftmark}}
\fancyhead[RO]{\small\color{forest}\nouppercase{\rightmark}\quad\thepage}
\renewcommand{\headrulewidth}{0.3pt}
\fancypagestyle{plain}{\fancyhf{}\fancyfoot[C]{\color{forest}\thepage}\renewcommand{\headrulewidth}{0pt}}

\titleformat{\chapter}[display]
  {\normalfont\bfseries\color{forest}}
  {\filleft\Large\chaptername\ \thechapter}
  {0.5ex}
  {\titlerule[1.2pt]\vspace{1.1ex}\Huge}
\titleformat{\section}{\Large\bfseries\color{forest}}{\thesection}{0.7em}{}
\titleformat{\subsection}{\large\bfseries\color{terracotta}}{\thesubsection}{0.7em}{}
\setlist{itemsep=0.25em,topsep=0.5em}

\newtheorem{theoreme}{Théorème}[chapter]
\newtheorem{proposition}[theoreme]{Proposition}
\newtheorem{lemme}[theoreme]{Lemme}
\newtheorem{corollary}[theoreme]{Corollaire}
\newtheorem{definition}[theoreme]{Définition}
\newtheorem{claim}[theoreme]{Affirmation}
\newtheorem{remark}[theoreme]{Remarque}
\newtheorem{exemple}{Exemple}[chapter]

\newenvironment{game}[3][]{%
  \renewcommand{\arraystretch}{1.15}%
  \setlength{\tabcolsep}{0.9em}%
  \begin{tabular}{r|*{#3}{c|}}
  \ifstrempty{#1}{}{\multicolumn{\numexpr#3+1\relax}{c}{\itshape #1}\\}}
  {\end{tabular}}

\newcommand\varal{\mathbf}
\newcommand\seq{\mathbf}
\newcommand\set{\mathcal}
\newcommand\EP{\mathbf{EP}}
\newcommand\E{\mathbf{E}}
\newcommand\AEP{\mathbf{AEP}}
\newcommand\indic{\mathbbm{1}_}
\newcommand\reals{\mathbb R}
\newcommand\naturals{\mathbb N}
\newcommand\TypeSet{\mathbb T}
\newcommand\spa{\,}
\let\vare\varepsilon
\let\Implies\implies
\let\implies\Rightarrow
\let\Iff\iff
\let\iff\Leftrightarrow
\newcommand\feq{\stackrel{\smash{\text{def}}}{=}}
\newcommand\restr[1]{_{\smash{\vert #1}}}
\def\p{\vskip 4pt \noindent}
\def\co{\operatorname{co}}
\def\cav{\operatorname{cav}}

\begin{document}
\frontmatter
\begin{titlepage}
\centering
\vspace*{0.16\textheight}
{\color{forest}\rule{\textwidth}{1.4pt}\par}
\vspace{1.6em}
{\Huge\bfseries Théorie de la décision\\[0.25em] et théorie des jeux\par}
\vspace{1.4em}
{\Large Notes de cours\par}
\vspace{2.2em}
{\LARGE Olivier Gossner\par}
\vfill
{\large Cours initialement enseigné à l'École nationale des ponts et chaussées\par}
\vspace{0.6em}
{\normalsize Édition remise en forme, 2026\par}
\vspace{1.6em}
{\color{terracotta}\rule{0.55\textwidth}{1pt}\par}
\end{titlepage}

\chapter*{Avant-propos}
Ces notes ont été préparées pour un cours de théorie de la décision et de théorie des jeux enseigné à l'École nationale des ponts et chaussées entre 2002 et 2004. Cette édition réunit les chapitres auparavant diffusés séparément, corrige diverses erreurs typographiques et en modernise la présentation. Le contenu scientifique reste celui du cours original.

Elles proposent une introduction aux préférences et à l'utilité espérée, à l'aversion au risque, à la décision dynamique, puis aux jeux sous forme normale et extensive, aux stratégies mixtes et aux principaux concepts d'équilibre.

\tableofcontents
\mainmatter
\chapter{Théorie de la d\'ecision}

\section{Introduction}
La théorie de la décision modélise le comportement d'un agent face
à des situations de choix. Nous décrivons les choix d'un agent, et
les relions à une notion de préférences. Nous montrons sous
quelles conditions les relations de pr\'ef\'erences peuvent \^etre
repr\'esent\'ees par une fonction d'utilit\'e, et quand elles
impliquent l'existence de croyances de l'agent sur les
diff\'erents \'etats du monde, repr\'esentables par des
probabilit\'es.

\begin{comment}Je crois qu'il serait bien de parler des deux aspects :
normatif et descriptif, car cela aide bien à comprendre le cours
par la suite, et est vrai pour toute la théorie des jeux et la
théorie économique en general.
\end{comment}

\section{Choix et préférences}
Nous commençons par décrire des préférences d'un agent, et des
choix, et établissons un lien naturel entre ces deux notions.
\subsection{Choix}
Nous partons d'un ensemble fini $X$ d'alternatives. Un agent peut
effectuer un choix parmi un sous-ensemble de $X$ d'alternatives
qui lui sont proposées.  Soit donc $2^X$ l'ensemble des parties
de~$X$.
\begin{definition}
Une fonction de choix est la donnée de $c\colon
2^X\backslash\{\emptyset\}\to 2^X\backslash\{\emptyset\}$ telle
que pour tout $A\subset X$, $c(A)\subset A$.
\end{definition}
L'interprétation est la suivante : si l'agent a le choix entre les
éléments de $A$, cet agent dira que tous les éléments de $c(A)$
conviennent, ou bien choisira indifféremment parmi les éléments de
$c(A)$.

 Certaines fonctions de choix vérifient des propriétés qui les
rendent plus ``plausibles'' que d'autres. Nous étudierons plus
particulièrement l'influence des axiomes suivant, appelés axiomes
$(\alpha)$ et $(\beta)$ de Sen :
\begin{itemize}
\item{\bf $(\alpha)$:} Si $x\in B\subset A$ et $x\in c(A)$, alors
$x\in c(B)$ %
\item{\bf $(\beta)$:} Si $x,y\in c(A)$, $A\subseteq B$, et $y\in
c(B)$, alors $x\in c(B)$.
\end{itemize}
On peut traduire ces axiomes en langage courant de la manière
suivante: Pour l'axiome $(\alpha)$, ``Si le champion du monde d'un
jeu est pakistanais, alors ce pakistanais doit aussi \^etre
champion du Pakistan.'' Pour l'axiome $(\beta)$, ``Si le champion
du monde est pakistanais, alors tous les champions du Pakistan
sont aussi champions du monde.''

\begin{comment}
En cours, donner des exemples de fonctions de choix, et montrer
qu'elles vérifient (ou non) certains des axiomes.
\end{comment}

\subsection{Préférences}

\begin{definition}
Une relation de préférences $\succ$ est la donnée d'un ensemble de
paires d'éléments $(x,y)$ de $X\times X$ pour lesquels nous
écrivons $x\succ y$. La relation $x\succ y$ est interprétée par :
``l'agent préf\`ere strictement $x$ à~$y$''.
\end{definition}

Voici une liste d'axiomes sur les relations de préférences:
\begin{itemize}
\item[\bf Asymétrie] $x\succ y\implies \neg y\succ x$. Si $x$ est
strictement préféré à $y$, alors $y$ n'est pas strictement préféré
à $x$. %
\item[\bf Transitivité] $(x\succ y\mbox{ et }y\succ z)\implies
x\succ z$. %
\item[\bf Irreflexivité] $\neg x\succ x$. Il n'existe pas de $x$
strictement préféré à lui-m\^eme.%
\item[\bf Transitivité négative] $(\neg x\succ y \mbox{ et }\neg
y\succ z)\implies \neg x\succ z$.%
\item[\bf Acyclicité] $x_1\succ x_2\succ\ldots\succ x_n\implies
\neg x_n\succ x_1$. Si chaque $x_i$ est strictement préféré à
$x_{i+1}$, alors $x_n$ n'est pas préféré à $x_1$.
\end{itemize}

\begin{definition}
Nous dirons qu'une relation de préférences est rationnelle si elle
est asymétrique et négativement transitive.
\end{definition}
Une relation de préférences ainsi définie possède forcément les
autres propriétés de la liste:
\begin{proposition}\label{rationelleimpliquelereste}
Une relation de préférences rationnelle est irreflexive,
transitive, et acyclique.
\end{proposition}
Une relation de préférences rationnelle est un ordre partiel.

A partir d'une relation binaire $\succ$, nous pouvons définir une
relation de préférences faibles $\succeq$  et une relation
d'indifférence $\sim$ par:
\begin{eqnarray*}
x\succeq y &\mbox{\ si\ }& \neg y\succ x \\
x\sim y &\mbox{\ si\ }& \neg y\succ x\mbox{\ et\ } \neg x\succ y %
\end{eqnarray*}

Voici quelques propriété des relations $\succeq$ et $\sim$ lorsque
$\succ$ est rationnelle.

\begin{proposition}
Supposons que $\succ$ est une relation de préférences rationnelle,
alors
\begin{itemize}
\item Pour tout $x$ et $y$, on a soit $x\succ y$, soit $y\succ x$,
soit $x\sim y$.%
\item $x\geq y$ est compl\^ete (on a toujours $x\geq y$ ou $y\geq
x$) et transitive. %
\item $\sim$ reflexive ($x\sim x$ pour tout $x$), symétrique
($x\sim y$ si et seulement si $y\sim x$), et transitive. %
\item $w\succ x$, $x\sim y$ et $y\succ z$ impliquent $w\succ y$ et
$x\succ z$.%
\item $x\geq y$ si et seulement si $x\succ y$ ou $x\sim y$. %
\item $x\geq y$ et $y\geq x$ si et seulement si $x\sim y$.
\end{itemize}
\end{proposition}

\subsection{Lien entre choix et préférences}
Deux idées simples et importantes relient les choix et les
préférences. Si un agent possède des préférences, ceci doit avoir
une implication en termes de choix. Lorsqu'on cherche quelles sont
les relations de préférences compatibles avec des choix (s'il en
existe), on parle de préférences révélées.

\subsubsection{Comment les préférences définissent des choix}
Soit $\succ$ une relation binaire sur $X$. On définit une fonction
$c(\cdot,\succ)$ par :
$$
c(A,\succ) = \{x\in A \colon \mbox{ pour tout $y\in A$ } \neq y \succ x\}
$$
$c(A,\succ)$ est donc le sous-ensemble des éléments de $A$ qui ne
sont pas dominés par d'autres éléments de $A$. Si $\succ$
représente les préférences de l'agents, alors ses choix doivent
obéir a $c(\cdot,\succ)$. Sous quelles conditions sommes-nous
s\^urs que $c(\cdot,\succ)$ est une fonction de choix ?

\begin{proposition}
Étant donnée une relation binaire $\succ$, $c(\cdot,\succ)$ est
une fonction de choix si et seulement si $\succ$ est acyclique.
\end{proposition}

Montrons que l'acyclicité implique $c(B,\succ)\neq\emptyset$.
C'est une consequence de $X$ fini. Réciproquement, $\succ$ n'est
pas acyclique, soit $x_1\succ x_2\succ\ldots\succ x_n\succ x_1$.
On a alors $c(\cdot, \succ)(\{x_1,x_2,\ldots, x_n\})=\emptyset$.

\subsubsection{Représentation des choix par des préférences rationnelles}
Nous partons maintenant d'une fonction de choix $c$, et cherchons
sous quelles conditions il existe une relation de préférences
rationnelle $\succ$ telle que $c$ soit la fonction de choix
correspondant à $\succ$ ($c(\cdot,\succ)=c$).

\begin{proposition}
Étant donnée une fonction de choix $c$, il existe une relation
rationnelle $\succ$ telle que $c=c(\cdot,\succ)$ si et seulement
si $\succ$ vérifie les axiomes $(\alpha)$ et $(\beta)$ de Sen. De
plus, cette relation est alors unique.
\end{proposition}

 Démonstration : Tout d'abord, supposons $c=c(\cdot,\succ)$, avec $\succ$
rationnelle. Pour montrer $\alpha$, il suffit de remarquer que
$\not\exists y \in A, y\succ x$ implique $\not\exists y \in B,
y\succ x$, pour $B$ sous-ensemble de $A$. Pour $\beta$, supposons
$(x,y)\in c(A)$, on a alors $x\sim y$. Si $x\in c(B)$ avec
$A\subset B$, $\forall z\in B$, $z\preceq x$, et donc $z\preceq y$
par transitivité, donc $y\in c(B)$.\\
 Supposons maintenant que $c$ vérifie $\alpha$ et $\beta$.
L'unicité de $\succ$ provient de ce que $\succ$ doit vérifier
$x\succ y$ si et seulement si $c(\{x,y\})= \{x\}$. La reflexivité
de $\succ$ est triviale. Pour la transitivité de $\succeq$,
supposons $x\succeq y\succeq z$. Si $y\in c(\{x,y,z\})$ alors
$z\in c(\{x,y,z\}$ ($\beta$). Sinon, on ne peut avoir
$c(\{x,y,z\}=\{x\}$ (car alors on aurait $c(\{x,y,z\}=\{x\}$).
Donc $z\in c(\{x,y,z\}=\{x\}$. On a donc aussi $z\in c(\{x,z\}$
($\alpha$), et $x\succeq z$. La relation $\succ$ est donc
rationnelle. Il reste à montrer que $c=c(\cdot, \succ)$, c'est à
dire que $x\in c(A)$ si et seulement si pour tout $y\in A$, $x\in
c(\{x,y\})$. Pour la partie ``seulement si'', $\forall y$, $x\in
c(\{x,y\})$, par l'axiome $\alpha$. Pour la partie ``si'', soit
$z\in c(A)$, alors $z\in c(\{x,z\})$ ($\alpha$), donc $c(\{x,z\})=
\{x,z\}$, et ($\beta$) implique $x\in c(A)$.


Il suffit d'observer les choix faits par les agents ($c$) pour en
déduire leurs préférences ($\succ$). Pour cette raison, la
conjonction des axiomes ($\alpha$) et ($\beta$) est aussi appelée
{\em axiome des préférences révélées}.

\section{Représentation ordinale}
Soit $u$ une fonction sur $X$. On interprète $u(x)$ comme une
utilité, ou plaisir, que l'agent peut retirer de l'alternative
$x$. On appelle donc $u$ {\em fonction d'utilité}. Si $u$ est la
fonction d'utilité de l'agent, quelles doivent \^etre ses
préférences ? Simplement, définissons $\succ_u$ par
$$
x\succ_u y \mbox{ si et seulement si } u(x) > u(y)
$$
Il est aisé de vérifier que $\succ_u$ est bien une relation de
préférences rationnelle. La proposition suivante nous dit aussi
que toute relation de préférences résulte d'une fonction d'utilité
($X$ est fini).
\begin{proposition}\label{repsiXfini}
Supposons que $X$ est fini, et soit $\succ$ une relation binaire
sur $X$. Il existe une fonction d'utilité $u$ telle que $\succ_u =
\succ$ si et seulement si $\succ$ est rationnelle.
\end{proposition}
Preuve : par induction sur le cardinal de $X$.

La proposition \ref{repsiXfini} s'étend au cas $X$ dénombrable.

\begin{comment}
Pourquoi ne pouvons pas faire d'induction transfinie ? Quelles
seraient les axiomes permettant une induction transfinie ?
Conjecture: probablement qqch + fort que l'acyclicité, genre pas
de cycles infinis...
\end{comment}

\begin{proposition}\label{repsiXdenom}
Supposons $X$ dénombrable, et soit $\succ$ une relation binaire
sur $X$. Il existe une fonction d'utilité $u$ telle que $\succ_u =
\succ$ si et seulement si $\succ$ est rationnelle.
\end{proposition}

Nous pouvons nous poser la question de l'unicité de la
représentation d'une relation de préférences par une fonction
d'utilité. Il est évident que la composition de $u$ par une
fonction strictement croissante ne change par $\succ_u$. La
proposition suivante nous donne une réciproque:

\begin{proposition}\label{unicitepreprordinale}
Soient $u$ et $u'$ deux fonctions d'utilité sur $X$ (fini). Alors
$\succ_u=\succ_{u'}$ si et seulement si il existe une fonction
$f\colon \reals \to \reals$ telle que:
\begin{itemize}
\item $f$ est strictement croissante%
\item $u'=f\circ u$
\end{itemize}
\end{proposition}
On commence par définir $f$ sur $\{u(x), x\in X\}$. Puis on
l'étend à $\reals$.

\section{Préférences sur choix incertains}
Introduction. Expliquer la différence entre choix incertains dont
on peut évaluer les probabilités de manière objective ou non.

\subsection{Préférences sur les probabilités objectives}

Soit $Z$ un ensemble fini, appelé ensemble de prix. Soit
$P=\Delta(Z)$ l'ensemble des lois de probabilités sur $Z$, appelé
ensemble des loteries. Nous étudions maintenant des relations de
préférences sur $P$. $P$  est donc l'ensemble de choix.

 Il est remarquable que $P$ a une structure particulière, appelée
structure de simplexe. Nommément :
\begin{itemize}
\item $P$ est convexe. Pour deux éléments $p$ et $q$ de $P$, et
$\lambda\in[0,1]$, $\lambda p + (1-\lambda) q \in P$, o\`u
$(\lambda p + (1-\lambda) q)(x)$ vaut par définition $\lambda p(x)
+ (1-\lambda) q(x)$. %
\item On peut identifier chaque $z\in Z$ à l'élément de $P$ qui
donne une probabilité de 1 à $z$, et $0$ aux autres.
\end{itemize}

 On définit maintenant des axiomes sur la relation de préférences
$\succ$ de $P$ qui utilisent la structure particulière de $P$.

\begin{itemize}
\item[{\bf Rationalité}] $\succ$ est une relation de
préférences rationnelle.%
\item[\bf Indépendance] Pour $p,q,r\in P$, et $0<\lambda\leq 1$,
$p\succ q$ implique $\lambda p+(1-\lambda)r\succ \lambda
q+(1-\lambda)r$%
\item[\bf Continuité] Pour $q\in P$, les ensembles $\{p, p\succ
q\}$ et $\{p, q\succ p\}$ sont ouverts
\end{itemize}

L'axiome de continuité se comprend lorsqu'on se rappelle que
$\succ$ représente une relation de préférences stricte. Si $p\succ
q$, alors si on modifie de manière extr\^emement mince $p$ en $p'$
et $q$ en $q'$ on aura $p'\succ q'$.

\subsection{Utilité espérée}
Supposons que l'agent associe une valeur, ou utilité $u(z)$ à
chacun des prix $z\in Z$. La fonction $u$ définit une fonction
$\tilde u$ de $P$ vers $\reals$ donnée par:
$$
\tilde u(p) = \sum_z p(z) u(z)
$$
Soit $\succ_{\tilde u}$ la relation de préférences sur $P$
associée à $\tilde u$.

Remarquons que nous n'avons aucune raison de supposer {\em a
priori} que les préférences de l'agent soient de la forme
$\succ_{\tilde u}$, pour un $u$ donné, ni encore qu'il existe $u$
telle que $\succ = \succ_{\tilde u}$. Il peut d'ailleurs sembler à
première vue que cette forme de préférences est restrictive et
très particulière. Cependant, cette forme de préférences est très
commode pour le modélisateur (nous), ou pour un agent cherchant à
évaluer ses préférences entre des loteries, car la donnée
d'utilités associées aux événements certains (les prix de $Z$) à
elle seule définit les préférences sur l'ensemble beaucoup plus
complexe des loteries.

\subsection{Théorème de von-Neumann et Morgenstern}
Le résultat de von-Neumann et Morgenstern est l'équivalence entre
une famille d'axiomes sur $\succ$ et l'existence d'une
représentation de la forme $\succ_{\tilde u}$.
%
\begin{proposition}
Il existe une fonction $u:Z\to \reals$ telle que $\succ =
\succ_{\tilde u}$ si et seulement si $\succ$ vérifie les axiomes
de préférences, d'indépendance et de continuité.
\end{proposition}
%
{\bf Démonstration de la partie ``si''} Il suffit de montrer que
$\succ_{\tilde u}$ vérifie les propriétés requises. \\
{\bf Démonstration de la partie ``seulement si''} Pour $z\in Z$,
on note $\delta_z$ l'élément de $P$ qui met probabilité 1 sur $Z$.
Si pour tous $x,y\in Z$, $\delta_x\sim\delta_y$, alors pour tous
$p,p'\in P$, $p\sim p'$ (indépendance), et toute fonction $u$
constante convient. Dans le cas contraire, on sélectionne parmi
l'ensemble $\{\delta_z,\ z\in Z\}$ un élément minimal
$\delta_{x_0}$, et un élément maximal $\delta_{x_1}$ pour $\succ$.
D'après l'axiome de continuité, pour chaque $z$, il existe un
unique $\lambda_z\in[0,1]$ tel que $z\sim\lambda_z
\delta_{x_1}+(1-\lambda_z)\delta_{x_0}$. On pose alors
$u(z)=\lambda_z$. On a donc $u(x_0)=0$, $u(x_1)=1$, et pour $p\in
P$ $\tilde u(p)= \sum_k p(z)\lambda(z)$. Soient $p$ et
$p'$ dans $P$.%
\begin{eqnarray*}
p\succ p' & \Leftrightarrow & \sum_z p'(z)(\lambda_z
\delta_{x_1}+(1-\lambda_z)\delta_{x_0}) \succ \sum_z
p'(z)(\lambda_z \delta_{x_1}+(1-\lambda_z)\delta_{x_0})\\
& \Leftrightarrow & (\sum_z p(z)\lambda(z))\delta_{x_1}+(\sum_z
p(z)(1-\lambda(z)))\delta_{x_0} \succ \\
&&(\sum_z p(z)'\lambda(z))\delta_{x_1}+(\sum_z
p(z)'(1-\lambda(z)))\delta_{x_0}\\
& \Leftrightarrow & \sum_z  p(z)\lambda(z) > \sum_z p'(z)\lambda(z)\\
& \Leftrightarrow & \tilde u(p) > \tilde u(p')
\end{eqnarray*}
\begin{proposition}
Deux fonctions d'utilité $\tilde u$ et $\tilde u'$ sur $P$ sont
associées aux m\^emes préférences si et seulement il existe $a>0$
et $b$ tels que $\tilde u' = a\tilde u+b$.
\end{proposition}
\noindent
{\bf Démonstration de la partie ``si''} est évidente.\\
{\bf Démonstration de la partie ``seulement si''} Supposons que
$\succ_{\tilde u} = \succ_{\tilde u'}$, et soient $\delta_{x_0}$,
$\delta_{x_1}$ un élément minimal et un élément maximal. Pour tout
$p\in P$, écrivons $\tilde u(p)= \frac{u(x_1)-\tilde
u(p)}{u(x_1)-u(x_0)}u(x_0)+\frac{\tilde
u(p)-u(x_0)}{u(x_1)-u(x_0)}u(x_1)$. Donc $p\sim_{\tilde
u}\frac{u(x_1)-\tilde
u(p)}{u(x_1)-u(x_0)}\delta_{x_0}+\frac{\tilde
u(p)-u(x_0)}{u(x_1)-u(x_0)}\delta_{x_1}$. Comme $\sim_{\tilde
u}=\sim_{\tilde u'}$,
$$%
\frac{u(x_1)-\tilde u(p)}{u(x_1)-u(x_0)} = \frac{u'(x_1)-\tilde
u(p)}{u'(x_1)-u'(x_0)}
$$%
Cad.
$$%
\tilde u'(p)= \frac{u'(x_1)-u'(x_0)}{u(x_1)-u(x_0)}u(p)-
\frac{u'(x_1)-u'(x_0)}{u(x_1)-u(x_0)}u(x_1) + u'(x_1)
$$%
ce qui est bien la forme recherchée.
\subsection{Préférences et incertain subjectif}
Dans une loterie, l'agent à un moyen objectif d'estimer la
probabilité de chaque prix. En revanche, dans un événement comme
une course de chevaux par exemple, il n'existe pas de moyen
objectif de conna\^itre la probabilité que tel ou tel cheval soit
vainqueur. La théorie de Anscombe et Aumann nous donne des
conditions sous lesquelles l'agent se comporte {\em comme si} il
possédait une croyance sur les chevaux vainqueurs, et effectuait
ses choix (ici des paris) en maximisant son utilité espérée étant
donnée cette croyance. Ces probabilités sont dites probabilités
subjectives.

\subsubsection{Paris sur des loteries}
$S$ est un ensemble fini de possibilités incertaines (comme les
chevaux gagnants dans le cas de la course). $Z$ est un ensemble
fini de prix, et $P$ représente les probabilités sur $Z$. A chaque
élément de $S$, faisons correspondre un prix aléatoire. Imaginons
que pour chaque élément de $S$, une loterie puisse \^etre tirée
pour déterminer le prix dans $Z$. On définit alors une {loterie
composée} comme un vecteur $h =(h_s)_{s\in S}$ de $P^S$.
L'ensemble de choix que nous considérons maintenant l'ensemble $H=
P^S$ des loteries composées.

Exemples de loteries composées :
\begin{itemize}
\item S'il pleut, mon cheval préféré Rossinante a 45\% de chances
de remporter la course, et seulement 20\% de chances s'il ne pleut
pas.%
\item Si l'économie est en croissance l'an prochain, j'ai une
chance sur deux d'obtenir une promotion, et une chance sur trois
sinon.
\end{itemize}

 Remarquons que $H$ a lui aussi une structure de simplexe. En
particulier, si $h,h'\in H$ et si $\lambda\in [0,1]$, $h''$ défini
par $h_s''(z)= \lambda h_s(z)+(1-\lambda)h_s'(z)$ appartient à
$H$.

\subsubsection{Les axiomes}
Supposons que l'agent a des préférences sur $H$, représentées par
$\succ$. Comme dans le cas des probabilités objectives, soient les
axiomes suivants sur $\succ$.

\begin{itemize}
\item[Axiome de rationalité] $\succ$ est une relation de
préférences rationnelle (sur $H$). \item[Axiome d'indépendance]
Pour $h,h',h''\in H$, et $0<\lambda\leq 1$, $h\succ h'$ implique
$\lambda h+(1-\lambda)h''\succ \lambda h+(1-\lambda)h'$
\item[Axiome de continuité] Pour $h,h',h''\in H$, si $h\succ
h'\succ h''$ il existe $0<\lambda,\mu<1$ tels que $\lambda
h+(1-\lambda)h''\succ h'\succ \mu h+(1-\mu)h''$
\end{itemize}

\subsubsection{Utilité dépendante de l'état}
Comme pour le cas des fonctions d'utilité de von-Neumann et
Morgenstern, nous pouvons caractériser les relations de binaires
qui vérifient les axiomes de  préférences, d'indépendance et de
continuité.
%
\begin{proposition}
Une relation binaire $\succ$ vérifie les axiomes de préférences,
d'indépendance et de continuité si et seulement si il existe une
famille $u=(u_s)_{s\in S}$ d'applications $u_s\colon Z\to \reals$
telles que
$$
h\succ h' \Leftrightarrow \sum_{s\in S}\sum_{z\in Z}
h_s(z)u_s(z)>\sum_{s\in S}\sum_{z\in Z} h'_s(z)u_s(z)
$$
\end{proposition}
\begin{comment}
Parler d'unicité ? En parler pour les fonctions de vN-M ?
\end{comment}
La démonstration est la m\^eme que celle du théorème de
von-Neumann et Morgenstern.
%
On peut alors poser $\tilde u_s(h_s)=\sum_z h_s(z)u_s(z)$, utilité
de von-Neumann et Morgenstern dans l'état du monde $s$.

\subsubsection{Utilité indépendante de l'état et probabilités subjectives}

 Nous allons introduire un axiome disant que les préférences ne
dépendent pas de l'état du monde. Soit donc $\succ_{\tilde u_s}$
la relation de préférences dans l'état $s$.
\begin{itemize}
\item[\bf Indépendance de l'état] $\succ_{\tilde u_s}$ ne dépend
pas de $s$.
\end{itemize}
Selon l'axiome d'indépendance de l'état, si $p$ est préféré à $p'$
dans un état $s$, il est aussi préféré à $p'$ dans tous les états
$s'$.

Dans ce cas, fixons un état $s_0$, et posons $\tilde u = \tilde
u_{s_0}$. On a pour tout état $s$ l'existence de réels
$a_{s}>0,b_s$ tels que , $\tilde u_{s}=
a_{s}\tilde u_{s_0} +b_{s}$. On peut donc écrire:%
$$
\sum_{s} \tilde u_s(h_s) = \sum_s a_s \tilde u_{s_0}(h_s) +\sum_s b_s
$$
 On a donc $h\succ h'$ si et seulement si:
$$
\sum_s a_s \tilde u_{s_0}(h_s) > \sum_s a_s \tilde u_{s_0}(h'_s)
$$
Posons alors $\pi(s) = \frac{a_s}{\sum_s a_s}$. Tous les $\pi(s)$
sont positifs, et leur somme vaut 1. On peut donc interpréter
$\pi=(\pi(s))_s$ comme une probabilité sur $S$. On a alors $h\succ
h'$ si et seulement si:
$$
\sum_s \pi(s) \tilde u_{s_0}(h_s) > \sum_s \pi(s) \tilde
u_{s_0}(h'_s)
$$
Tout se passe donc comme si les préférences de l'agent étaient
représentées par les préférences $\tilde u$, et par une
probabilité $\pi$ sur $S$. Il est naturel d'interpréter $\pi$
comme les {\em croyances} de l'agent sur $S$. Ces croyances ne
sont pas une donnée objective, on parle de croyances {\em
subjectives} de l'agent.
%
\section{Exercices}
\subsection{Exercice 1}
Soit l'axiome suivant, appelé Axiome de Houthakker, portant sur
une fonction de choix $c$:\\
Si $x,y\in A\cap B$, et si $x\in c(A)$ et $y\in c(B)$, alors $x\in
c(B)$.\\
Montrer que l'axiome de Houthakker est équivalent à la conjonction
des axiomes $\alpha$ et $\beta$ de Sen.
\subsection{Exercice 2}
Montrer la proposition \ref{rationelleimpliquelereste} %
\subsection{Exercice 3}
Une relation de préférences non représentable par une fonction
d'utilité. Sur $X=[0,1]\times[0,1]$, on définit la relation de
préférences entre $x=(x_1,x_2)$ et $y=(y_1,y_2)$ par $x\succ y$ si
et seulement si $x_1> y_1$, ou ($x_1=y_1$ et $x_2>y_2$). C'est
donc l'ordre lexicographique sur $X$.
\begin{enumerate}
\item Montrer que $\succ$ est rationnelle.%
\item Montrer que $\succ$ n'est représentable par aucune fonction
d'utilité. C'est à dire, montrer qu'il n'existe pas d'application
$u$ de $X$ vers $\reals$ telle que $x\succ y$ si et seulement si
$u(x)> u(y)$.
\end{enumerate}
\subsection{Exercice 5} On a deux actifs financiers. Le premier
paye 100 dollars dans l'état du monde 1 (reprise de l'économie
américaine), et 50 dollars dans l'état du monde 2 (pas de reprise
de l'économie américaine). Le deuxième paye 120 dollars dans
l'état du monde 1 et 20 dans l'état du monde 2. Sur les marchés,
le prix du premier actif financier est de $80$ euros, celui du
second de 70 euros. Dans le ``Financial Times'' vous lisez le
commentaire suivant, ``d'après le marché, l'état du monde 1 et
l'état du monde 2 ont les mêmes chances de se réaliser.''
\begin{enumerate}
\item Ce commentaire est-il en rapport avec les prix que vous
observez ? %
\item  Selon vous, quelles sont les probabilités subjectives du
marché sur l'état 1 et sur l'état 2 ? %
\end{enumerate}
\subsection{Exercice 6: Paradoxe d'Allais}
On a comme ensemble d'alternatives $X =\{ 2 500 000 $\euro$, 500
000 $\euro$, 0 $\euro$ \}$, où $2 500 000 $\euro correspond au
fait de gagner cette somme, etc.\ldots On a les loteries $p_1 =
(0,1,0)$,  $p'_1 = (0.1,0.89,0.01)$, $p_2 = (0,0.11,0.89)$ et $p'2
= (0.1,0,0.9)$. (Par exemple, avec $p_2$, la probabilité de gagner
2500000\euro~ est de 0.11.
\begin{enumerate}
\item On suppose qu'un individu préfère $p_1$ à $p'_1$, et $p'_2$ à $p_2$.
Est-ce consistant avec l'utilité espérée de von-Neumann et Morgenstern ?
\end{enumerate}
\subsection{Exercice 7: Paradoxe de Machina}
$X = \{ $voyage à Venise, beau film sur Venise, rien $\}$, et $p_1
= (0.999,0.001,0)$, $p'_1 = (0.999,0,0.001)$. Si un agent préfère
voir un beau film sur Venise plut\^ot que rien, et préfère $p'_1$
à $p_1$, est-ce consistant avec la théorie de l'utilité espérée de
von-Neumann et Morgenstern ? Comment peut-on expliquer de telles
préférences ?

\chapter{Aversion au risque}

La théorie de l'utilité de von Neumann et Morgenstern nous dit que
l'utilité associée à une loterie donnant un gain monétaire $x$
avec probabilité $\frac12$ et $y$ avec probabilité $\frac12$ est
la moyenne entre les utilités de $x$ et de $y$. Ce que la théorie
de von Neumann et Morgenstern ne nous dit pas, c'est comment se
comparent cette utilité avec celle du gain monétaire de $\frac12 x
+ \frac12 y$. C'est ce que nous étudions dans ce chapitre, qui
concerne les préférences de l'agent face au risque.

Les objectifs de ce chapitre sont :%
\begin{itemize}
\item D\'efinir un agent averse au risque ou qui aime le risque
\item Caractériser les fonctions d'utilité correspondantes. %
\item Étudier l'évolution du comportement face au risque avec la
richesse.
\item Mesurer cette aversion au risque ou cet amour du risque %
\end{itemize}

\section{Risque et utilités}

Les gains et les pertes sont mesurés de manière monétaire. On
considère donc des loteries à valeurs dans $\reals$, et soit $P=
\{$ loteries sur $\reals$ \`a support fini $\}$. En particulier,
pour $z\in \reals$, $\delta_z\in P$ représente la loterie qui
donne $z$ avec probabilité 1. Pour $z\in P$, $e(z)=\sum_x z(x).x$
est l'espérance de $p$. Pour $f\colon \reals\to\reals$, et $p\in
P$, $\E_p f$ est l'esp\'erance de $f$ sous $p$.

On se dote d'une relation de pr\'ef\'erences rationnelle $\succ$
sur $P$ repr\'esent\'ee par une fonction d'utilit\'e de von
Neumann--Morgenstern $u$.

On caractérise d'abord la croissance de l'utilité avec l'argent.
\begin{proposition}
$u$ est strictement croissante si et seulement si:
$$
\delta_z \succ \delta_{z'} \Leftrightarrow z > z'
$$
\end{proposition}

On supposera par la suite $u$ strictement croissante.

\begin{definition}~
\begin{itemize}
\item[$\bullet$] $\succ$ est {\bf averse au risque} si
$\delta_{e(p)}\succcurlyeq p$ pour tout $p\in P$.%
\item[$\bullet$] $\succ$ est {\bf strictement averse au risque} si
$\delta_{e(p)}\succ p$ pour tout $p\in P$.%
\item[$\bullet$] $\succ$ est {\bf neutre au risque} si
$\delta_{e(p)}\sim p$ pour tout $p\in P$.%
\item[$\bullet$] $\succ$ {\bf aime le risque} si
$\delta_{e(p)}\preccurlyeq p$ pour tout $p\in P$ %
\item[$\bullet$] $\succ$ {\bf aime strictement le risque} si
$\delta_{e(p)}\prec p$ pour tout $p\in P$.
\end{itemize}
\end{definition}
On a les caractérisations suivantes en termes de fonctions
d'utilité.
\begin{proposition}~
\begin{itemize}
\item[$\bullet$] $\succ$ est averse au risque si et seulement si
$u$ est
concave.%
\item[$\bullet$] $\succ$ est strictement averse au risque si et
seulement si
$u$ est strictement concave.%
\item[$\bullet$] $\succ$ est neutre au risque si $u$ est affine. %
\item[$\bullet$] $\succ$ aime le risque si $u$ est convexe. %
\item[$\bullet$] $\succ$ aime strictement le risque si $u$ est
strictement convexe.
\end{itemize}
\end{proposition}

\section{Risque et richesse}
Un agent aimera-t-il plus ou moins le risque lorsque sa richess
augmente ? Deux d\'efinitions sont possibles :
\begin{itemize}
\item Un m\^eme risque (en valeurs {\em absolues}) sera-t-il pris
par un agent plus riche, moins riche ?%
\item Un m\^eme risque (en valeurs {\em relatives}) sera-t-il pris
par un agent plus riche, moins riche ?
\end{itemize}

\subsection{Risque absolu}
On définit l'aversion au risque absolu en termes de préférences
avant de la caractériser en termes d'utilités.

\begin{definition}
$u$ (ou $\succ$) pr\'esente une {\em aversion au risque (absolu)
d\'ecroissante} [resp.\, croissante, resp.\, constante] si
$\forall$ $q\in P$, $z,w,w' \in \reals$ t.q.\, $w'>w$:
\begin{eqnarray*}
\E_{w+q}u>u(w+z) &\Rightarrow& \E_{w'+q}u>u(w'+z)\\
\mbox{resp.\,} &\Leftarrow& \\
\mbox{resp.\,} &\Leftrightarrow&
\end{eqnarray*}
\end{definition}

Dans le premier cas, si la loterie $q$ est pr\'ef\'er\'ee \`a la
somme $z$ s\^ure lorsque la richesse est $w$, elle l'est aussi si
la richesse est $w'>w$.


\begin{proposition} Supposons $u$ deux fois diff\'erentiable.
Alors $u$ pr\'esente une aversion au risque d\'ecroissante
[resp.\, croissante, resp.\, constante] si et seulement si la
fonction:
$$
\lambda = - \frac{u''}{u'}
$$
est d\'ecroissante [resp.\, d\'ecroissante, resp.\, constante]. \\
\end{proposition}

$\lambda$ s'appelle le {\bf coefficient d'aversion au risque de
Arrow-Pratt}. Il s'agit de la courbure de $u$ normalis\'ee par la
d\'eriv\'ee.


La proposition suivante caractérise l'aversion pour le risque
absolu constante:

\begin{proposition} Supposons $u$ deux fois diff\'erentiable. Alors
$u$ pr\'esente une aversion au risque constante si et seulement si
il existe $a>0$ et $b$ tels que:
$$
u(z)=
\begin{cases}
  az+b &\mbox{si }\;\; \lambda(z)\equiv0\cr
  -ae^{-\lambda z}+b  &\mbox{si }\;\; \lambda(z)\equiv\lambda>0
\end{cases}
$$
\end{proposition}

C'est-à-dire que les préférences sont représentables par une
fonction d'utilité $u$ qui peut être de deux types :
$$
u(z)=
\begin{cases}
  z &\mbox{si }\;\; \lambda(z)\equiv0\cr
  -e^{-\lambda z}  &\mbox{si }\;\; \lambda(z)\equiv\lambda>0
\end{cases}
$$

\subsection{Risque relatif}

\begin{definition}
$u$ (ou $\succ$) pr\'esente une {\em aversion au risque (relatif)
d\'ecroissante} [resp.\, croissante, resp.\, constante] si
$\forall$ $q\in P$ t.q.\ $q\geq 0$ avec pba.\, 1, $z>0$ $w'>w$:
\begin{eqnarray*}
\E_{wq}u>u(wz) &\Rightarrow& \E_{w'q}u>u(w'z)\\
\mbox{resp.\,} &\Leftarrow& \\
\mbox{resp.\,} &\Leftrightarrow&
\end{eqnarray*}
\end{definition}

Dans le premier cas, si on pr\'ef\`ere multiplier sa richesse par
la loterie $q$ plut\^ot que par $z$ lorsque la richesse est $w$,
on le pr\'ef\`ere aussi si la richesse est $w'>w$.

\begin{proposition} Supposons $u$ deux fois diff\'erentiable. Alors
$u$ pr\'esente une aversion au risque relatif d\'ecroissante
[resp.\, constante] si et seulement si la fonction:
$$
z\to  - \frac{zu''(z)}{u'(z)}
$$
est d\'ecroissante [resp.\, constante].
\end{proposition}

\begin{proposition} Supposons $u$ deux fois diff\'erentiable.
Alors $u$ pr\'esente une aversion au risque relatif constante si
et seulement si il existe $a>0$ et $b$ tels que:
$$
u(z)=
\begin{cases}
  a\ln(z)+b &\cr
  a \gamma z^{\gamma}+b  &\mbox{pour un certain }\gamma<1, \gamma\neq 0
\end{cases}
$$
\end{proposition}

C'est-à-dire que les préférences sont représentables par une
fonction d'utilité $u$ qui est d'un des trois types suivants :
$$
u(z)=
\begin{cases}
  \ln(z) & \cr
  - z^{\gamma}  &\mbox{pour un certain }\gamma<0\cr
    z^{\gamma}  &\mbox{pour un certain } 0<\gamma<1
\end{cases}
$$

\section{Application}
Supposons que vous ayez le choix entre les loteries suivantes:\\
\begin{minipage}{3.5cm}%
\begin{tabular}{p{0.4cm}p{0.4cm}p{0.4cm}}
&\rnode{a}{$\blue\bullet$}&\\[1.5cm]
\rnode{b}{$\blue\bullet$}&\rnode{c}{$\blue\bullet$}&\rnode{d}{$\blue\bullet$}\\[-0.4cm]
\red\tiny25\euro&\red\tiny150\euro&\red\tiny600\euro\\
\ncline[linewidth=1.3pt,linecolor=yellow]{a}{b}\aput[2pt]{0}(0.89){\tiny$0.7$}
\ncline[linewidth=1.3pt,linecolor=yellow]{a}{c}\aput[2pt]{0}(0.95){\tiny$0.2$}
\ncline[linewidth=1.3pt,linecolor=yellow]{a}{d}\aput[2pt]{0}(0.99){\tiny$0.1$}
\end{tabular}%
\end{minipage}%
\begin{minipage}{3.5cm}%
\begin{tabular}{p{0.4cm}p{0.4cm}p{0.4cm}}
&\rnode{a}{$\blue\bullet$}&\\[1.5cm]
\rnode{b}{$\blue\bullet$}&\rnode{c}{$\blue\bullet$}&\rnode{d}{$\blue\bullet$}\\[-0.4cm]
\red\tiny80\euro&\red\tiny90\euro&\red\tiny98\euro\\
\ncline[linewidth=1.3pt,linecolor=yellow]{a}{b}\aput[2pt]{0}(0.89){\tiny$0.2$}
\ncline[linewidth=1.3pt,linecolor=yellow]{a}{c}\aput[2pt]{0}(0.95){\tiny$0.58$}
\ncline[linewidth=1.3pt,linecolor=yellow]{a}{d}\aput[2pt]{0}(0.99){\tiny$0.22$}
\end{tabular}%
\end{minipage}%
\begin{minipage}{3.5cm}%
\begin{tabular}{p{0.4cm}p{0.4cm}p{0.4cm}}
&\rnode{a}{$\blue\bullet$}&\\[1.5cm]
\rnode{b}{$\blue\bullet$}&\rnode{c}{}&\rnode{d}{$\blue\bullet$}\\[-0.4cm]
\red\tiny2\euro&\red\tiny&\red\tiny105\euro\\
\ncline[linewidth=1.3pt,linecolor=yellow]{a}{b}\aput[2pt]{0}(0.89){\tiny$0.05$}
%\ncline[linewidth=1.3pt,linecolor=yellow]{a}{c}\aput[2pt]{0}(0.95){\tiny$0.58$}
\ncline[linewidth=1.3pt,linecolor=yellow]{a}{d}\aput[2pt]{0}(0.99){\tiny$0.95$}
\end{tabular}%
\end{minipage}\\
A-t-on une mani\`ere objective de choisir entre ces loteries~?

Une m\'ethode possible est de supposer une fonction d'utilité de
von Neumann et Morgenstern, avec une aversion au risque constante
{\em pour ces valeurs mon\'etaires}. Il reste \`a déterminer le
coefficient d'aversion au risque $\lambda$.

On peut déterminer $\lambda$ en répondant à la question suivante :
Combien donneriez-vous pour une loterie offrant $0$ avec
probabilit\'e $\frac12$ et $500$\,\euro\ avec probabilit\'e
$\frac12$ ?

Il est ensuite possible d'évaluer l'utilité pour chacune des
loteries, et de déterminer celle qui maximise l'utilité.

\chapter{Théorie de la décision dynamique}

\section{Introduction}
Nous considérons des situations dans lesquelles une suite de
décisions doit \^etre prise par l'agent. On montre comment décrire
ces situations par des arbres de décision, et comment décrire un
comportement de l'agent par une stratégie. On verra aussi comment
rechercher les stratégies optimales de l'agent, c'est-à-dire
celles qui maximisent son utilité.

\section{Décision dynamique en information parfaite}

\subsection{Description d'un problème de décision}
Un probl\`eme de décision dynamique en information parfaite est
donné par
\begin{itemize}
\item $T$ ensemble fini de n{\oe}uds, $Z\subset T$ n{\oe}uds
terminaux, et pour $t\not\in Z$ :
\begin{itemize}
\item  $i(t)\in \{N,J\}$ qui détermine si la nature ($N$) ou le
joueur ($J$) joue en $t$ ;%
\item L'ensemble de choix possibles au n{\oe}ud $t$, $A(t)$ ; %
\item Le n{\oe}ud successeur r\'esultant du choix $a$, $N(t,a)$ ;
\item Une distribution $D(t)$ sur $A(t)$ si $i(t)= N$.
\end{itemize}
\item $u:Z\to\reals$ fonctions d'utilit\'e du joueur ;
\end{itemize}

On définit ainsi les successeurs d'un n{\oe}ud
\begin{itemize}
\item $s(t)=\{N(t,a),a\in A(t)\}$ ensemble des successeurs directs
de $t$. ($\emptyset$ si $t\in Z$).%
\item $S(t)= s(t)\cup s(s(t)) \cup s(s(s(t)))...$ tous les
successeurs de $t$.
\end{itemize}

On suppose que le problème a une structure d'arbre, c'est à dire :

\begin{itemize}
\item Il n'y a pas de cycle : $\forall t, t\not\in S(t)$ %
\item Il existe un  n{\oe}ud initial $\exists\tilde t, S(\tilde
t)\cup\{\tilde t\}=T$ %
\item Chaque n{\oe}ud a un unique prédécesseur
\end{itemize}

\begin{exemple}
Exemple : casino.\\
$\bullet$ 1\`ere \'etape : Le joueur peut d\'ecider de jouer ou non.\\
S'il ne joue pas le jeu s'arr\^ete.\\
Sinon, la nature tire \`a pile ou face s'il gagne 1 ou perd 2.\\[0.5cm]
$\bullet$ 2\`eme \'etape : Le joueur peut d\'ecider de jouer ou non.\\
S'il ne joue pas le jeu s'arr\^ete.\\
S'il joue, la nature tire \`a pile ou face s'il gagne 3 ou perd 2
(en plus des gains et pertes de la première étape).

On représente ce jeu par l'arbre suivant :\\[0.5cm]
\begin{tabular}{p{1cm}p{1cm}p{1cm}p{1cm}p{1cm}p{1cm}}
&&&\rnode{a}{$\bullet$}&\rnode{c}{$0$}\\[0.7cm]
&&\rnode{b}{$\blue\diamond$}&&&\\[0.7cm]
&\rnode{d}{$\bullet$}&&\rnode{e}{$\bullet$}&\\[0.7cm]
&\rnode{f}{$\blue\diamond$}&\rnode{g}{\small$+1$}&\rnode{h}{$\blue\diamond$}&\rnode{i}{\small$-2$}&\\[0.7cm]
\rnode{j}{\small$+4$}&\rnode{k}{\small$-1$}&\rnode{l}{\small$+1$}&\rnode{m}{\small$-4$}\\[0.7cm]
\ncline[linewidth=0.7pt,linecolor=yellow]{-}{a}{b}\bput[1pt]{0}(0.5){\small$J_1$}%
\ncline[linewidth=0.7pt,linecolor=yellow]{-}{a}{c}\aput[1pt]{0}(0.5){\small$S_1$}%
\ncline[linewidth=0.7pt,linecolor=yellow]{-}{b}{d}\bput[1pt]{0}(0.5){\small$\frac12$}%
\ncline[linewidth=0.7pt,linecolor=yellow]{-}{b}{e}\aput[1pt]{0}(0.5){\small$\frac12$}%
\ncline[linewidth=0.7pt,linecolor=yellow]{-}{d}{f}\bput[1pt]{0}(0.5){\small$J_2$}%
\ncline[linewidth=0.7pt,linecolor=yellow]{-}{d}{g}\aput[1pt]{0}(0.5){\small$S_2$}%
\ncline[linewidth=0.7pt,linecolor=yellow]{-}{e}{h}\bput[1pt]{0}(0.5){\small$J_3$}%
\ncline[linewidth=0.7pt,linecolor=yellow]{-}{e}{i}\aput[1pt]{0}(0.5){\small$S_3$}%
\ncline[linewidth=0.7pt,linecolor=yellow]{-}{f}{j}\bput[1pt]{0}(0.5){\small$\frac12$}%
\ncline[linewidth=0.7pt,linecolor=yellow]{-}{f}{k}\aput[1pt]{0}(0.5){\small$\frac12$}%
\ncline[linewidth=0.7pt,linecolor=yellow]{-}{h}{l}\bput[1pt]{0}(0.5){\small$\frac12$}%
\ncline[linewidth=0.7pt,linecolor=yellow]{-}{h}{m}\aput[1pt]{0}(0.5){\small$\frac12$}%
%&\circlenode*{b}{\makebox[0.5cm]{{\red$B^{\blue11}$}}}&\circlenode*{e}{\makebox[0.5cm]{{\red$E^{\blue4}$}}}&&\\[0.1cm]
%&&&\circlenode*{h}{\makebox[0.5cm]{{\red$H^{\blue3}$}}}&\\[0.1cm]
%\circlenode*{a}{\makebox[0.5cm]{{\red$A^{\blue11}$}}}&\circlenode*{c}{\makebox[0.5cm]{{\red$C^{\blue7}$}}}%
\end{tabular}\\
Les n{\oe}uds $\bullet$ sont ceux où le joueur prend une décision.
Les décisions possibles au premier  n{\oe}ud de l'agent sont $J_1$
et $S_1$, $J_2$ et $S_2$ pour le deuxième, et $J_3$ et $S_3$ pour
le troisième. Les n{\oe}uds $\diamond$ sont ceux où la nature tire
au hasard. Les branches suivant ces n{\oe}uds sont affectées des
probabilités que la nature accorde à chacun d'entre eux. Les
n{\oe}uds affectés de nombres sont les n{\oe}uds terminaux, où
nous avons représenté le gain monétaire final de l'agent.

\end{exemple}

\subsection{Description d'un comportement : stratégie}

Une stratégie $s$ pour l'agent associe une décision dans $A(t)$ à
tout n{\oe}ud $t$ tel que $i(t)=J$. C'est donc une application de
$i^{-1}(J)$ vers $\cup_t A(t)$ telle que $s(t)\in A(t)$ pour tout
$t$.

A toute stratégie du joueur correspond une probabilité $p(s)$ sur
les n{\oe}uds terminaux, éléments de $Z$. Pour calculer $p(s)$, on
part du n{\oe}ud initial, et on descend progressivement d'une
étape à la fois, pour calculer les probabilités des successeurs,
en utilisant à un n{\oe}ud $t$ tel que $i(t)=N$ la probabilité
$D(t)$, et à un n{\oe}ud $t$ tel que $i(t)=J$ la décision $s(t)$.

On a donc une utilité espérée $U(s) = \E_{p(s)} u(z)$ pour chaque
stratégie. Une stratégie maximisatrice est une stratégie qui
maximise ce paiement espéré.

\begin{exemple}
On a dans le problème précédent 3 n{\oe}uds de décision pour
l'agent, avec à chacun d'entre eux 2 décisions possibles. Par
conséquent il y a 8 stratégies possibles. $J_1S_2S_3$ joue à la
première étape et s'arr\^ete ensuite. $J_1S_2J_3$ joue a la
première étape, puis continue si elle a perdu à la première, et
arr\^ete sinon. $S_1J_2J_3$ peut s'entendre comme ``ne pas jouer,
mais si on avait joué à la première étape, on aurait décidé de
continuer à la seconde''. Ce n'est pas la m\^eme chose que
$S_1S_2S_3$ qui d'arr\^ete à la première étape, mais décide aussi
de s'arr\^eter si on lui demande ce qu'elle ferait à la deuxième.

Les gains espérés de ces stratégies sont :
\begin{itemize}
\item $J_1S_2S_3$ donne $-\frac12$ %
\item $J_1S_2J_3$ donne $0$%
\item $S_1J_2J_3$ et $S_1S_2S_3$ donnent $0$.
\end{itemize}
\end{exemple}
%
\subsection{Principe d'induction à rebours}

Le principe d'induction à rebours permet de trouver une stratégie
maximisatrice. On part du bas de l'arbre, et pour les dernières
décisions de l'agent, on choisit la meilleure action. On remplace
ensuite le n{\oe}ud correspondant par l'utilité espérée étant
donnée cette meilleure action. On obtient ainsi un nouveau
problème de décision dynamique, auquel on applique de nouveau le
principe, et ainsi de suite... A la fin de la procédure, on aura
choisi une décision pour chaque n{\oe}ud, on aura donc obtenu une
stratégie dans le problème de décision. On a le résultat suivant :

\begin{proposition}
Toute stratégie obtenue par le principe d'induction à rebours est
maximisatrice.
\end{proposition}

\begin{exemple}
Question : Comment maximiser la probabilit\'e de finir avec un
gain d'au moins 1 dans l'exemple précédent ? Ceci revient à
chercher une stratégie maximisatrice lorsque l'utilité vaut 1 pour
les gains de 1 au moins, et 0 sinon. On a alors l'arbre :\\[0.5cm]
\begin{tabular}{p{1cm}p{1cm}p{1cm}p{1cm}p{1cm}p{1cm}}
&&&\rnode{a}{$\bullet$}&\rnode{c}{$0$}\\[0.7cm]
&&\rnode{b}{$\blue\diamond$}&&&\\[0.7cm]
&\rnode{d}{$\bullet$}&&\rnode{e}{$\bullet$}&\\[0.7cm]
&\rnode{f}{$\blue\diamond$}&\rnode{g}{\small$+1$}&\rnode{h}{$\blue\diamond$}&\rnode{i}{\small$0$}&\\[0.7cm]
\rnode{j}{\small$1$}&\rnode{k}{\small$0$}&\rnode{l}{\small$1$}&\rnode{m}{\small$0$}\\[0.7cm]
\ncline[linewidth=0.7pt,linecolor=yellow]{-}{a}{b}\bput[1pt]{0}(0.5){\small$J_1$}%
\ncline[linewidth=0.7pt,linecolor=yellow]{-}{a}{c}\aput[1pt]{0}(0.5){\small$S_1$}%
\ncline[linewidth=0.7pt,linecolor=yellow]{-}{b}{d}\bput[1pt]{0}(0.5){\small$\frac12$}%
\ncline[linewidth=0.7pt,linecolor=yellow]{-}{b}{e}\aput[1pt]{0}(0.5){\small$\frac12$}%
\ncline[linewidth=0.7pt,linecolor=yellow]{-}{d}{f}\bput[1pt]{0}(0.5){\small$J_2$}%
\ncline[linewidth=0.7pt,linecolor=yellow]{-}{d}{g}\aput[1pt]{0}(0.5){\small$S_2$}%
\ncline[linewidth=0.7pt,linecolor=yellow]{-}{e}{h}\bput[1pt]{0}(0.5){\small$J_3$}%
\ncline[linewidth=0.7pt,linecolor=yellow]{-}{e}{i}\aput[1pt]{0}(0.5){\small$S_3$}%
\ncline[linewidth=0.7pt,linecolor=yellow]{-}{f}{j}\bput[1pt]{0}(0.5){\small$\frac12$}%
\ncline[linewidth=0.7pt,linecolor=yellow]{-}{f}{k}\aput[1pt]{0}(0.5){\small$\frac12$}%
\ncline[linewidth=0.7pt,linecolor=yellow]{-}{h}{l}\bput[1pt]{0}(0.5){\small$\frac12$}%
\ncline[linewidth=0.7pt,linecolor=yellow]{-}{h}{m}\aput[1pt]{0}(0.5){\small$\frac12$}%
%&\circlenode*{b}{\makebox[0.5cm]{{\red$B^{\blue11}$}}}&\circlenode*{e}{\makebox[0.5cm]{{\red$E^{\blue4}$}}}&&\\[0.1cm]
%&&&\circlenode*{h}{\makebox[0.5cm]{{\red$H^{\blue3}$}}}&\\[0.1cm]
%\circlenode*{a}{\makebox[0.5cm]{{\red$A^{\blue11}$}}}&\circlenode*{c}{\makebox[0.5cm]{{\red$C^{\blue7}$}}}%
\end{tabular}

Appliquons le principe d'induction à rebours à cette arbre. On
voit qu'il faut décider $S_2$ au deuxième n{\oe}ud, et $J_3$ au
troisième. On obtient l'arbre dans lequel on a remplacé les
derniers n{\oe}uds de décisions par le gain espéré maximal :\\

\begin{tabular}{p{1cm}p{1cm}p{1cm}p{1cm}p{1cm}p{1cm}}
&&&\rnode{a}{$\bullet$}&\rnode{c}{$0$}\\[0.7cm]
&&\rnode{b}{$\blue\diamond$}&&&\\[0.7cm]
&\rnode{d}{$1$}&&\rnode{e}{$\frac12$}&\\[0.7cm]
\ncline[linewidth=0.7pt,linecolor=yellow]{-}{a}{b}\bput[1pt]{0}(0.5){\small$J_1$}%
\ncline[linewidth=0.7pt,linecolor=yellow]{-}{a}{c}\aput[1pt]{0}(0.5){\small$S_2$}%
\ncline[linewidth=0.7pt,linecolor=yellow]{-}{b}{d}\bput[1pt]{0}(0.5){\small$\frac12$}%
\ncline[linewidth=0.7pt,linecolor=yellow]{-}{b}{e}\aput[1pt]{0}(0.5){\small$\frac12$}%
%&\circlenode*{b}{\makebox[0.5cm]{{\red$B^{\blue11}$}}}&\circlenode*{e}{\makebox[0.5cm]{{\red$E^{\blue4}$}}}&&\\[0.1cm]
%&&&\circlenode*{h}{\makebox[0.5cm]{{\red$H^{\blue3}$}}}&\\[0.1cm]
%\circlenode*{a}{\makebox[0.5cm]{{\red$A^{\blue11}$}}}&\circlenode*{c}{\makebox[0.5cm]{{\red$C^{\blue7}$}}}%
\end{tabular}

On voit que dans ce jeu la meilleure décision est de jouer $J_1$
au premier n{\oe}ud. On a donc obtenu comme stratégie la stratégie
$J_1S_2J_3$ qui maximise la probabilité de terminer le jeu avec un
paiement d'au moins~1.
\end{exemple}

Une remarque importante sur la notion de stratégie et sur le
principe d'induction à rebours. Le principe d'action à rebours
indique quelle choix prendre, même aux n{\oe}uds qui ne devraient
pas être atteints. Par exemple, regardons le jeu suivant~:\\[0.5cm]

\begin{tabular}{p{1cm}p{1cm}p{1cm}p{1cm}p{1cm}p{1cm}}
&&&\rnode{a}{$\bullet$}&\rnode{c}{$0$}\\[0.7cm]
&&\rnode{b}{$\bullet$}&&&\\[0.7cm]
&\rnode{d}{$-1$}&&\rnode{e}{$-2$}&\\[0.7cm]
\ncline[linewidth=0.7pt,linecolor=yellow]{-}{a}{b}\bput[1pt]{0}(0.5){\small$J_1$}%
\ncline[linewidth=0.7pt,linecolor=yellow]{-}{a}{c}\aput[1pt]{0}(0.5){\small$S_1$}%
\ncline[linewidth=0.7pt,linecolor=yellow]{-}{b}{d}\bput[1pt]{0}(0.5){\small$J_2$}%
\ncline[linewidth=0.7pt,linecolor=yellow]{-}{b}{e}\aput[1pt]{0}(0.5){\small$S_2$}%
%&\circlenode*{b}{\makebox[0.5cm]{{\red$B^{\blue11}$}}}&\circlenode*{e}{\makebox[0.5cm]{{\red$E^{\blue4}$}}}&&\\[0.1cm]
%&&&\circlenode*{h}{\makebox[0.5cm]{{\red$H^{\blue3}$}}}&\\[0.1cm]
%\circlenode*{a}{\makebox[0.5cm]{{\red$A^{\blue11}$}}}&\circlenode*{c}{\makebox[0.5cm]{{\red$C^{\blue7}$}}}%
\end{tabular}

Le principe d'induction à rebours nous dit tout d'abord qu'il faut
choisir $J_2$ au deuxième n{\oe}ud, et transforme le problème en :\\[0.5cm]

\begin{tabular}{p{1cm}p{1cm}p{1cm}p{1cm}p{1cm}p{1cm}}
&&&\rnode{a}{$\bullet$}&\rnode{c}{$0$}\\[0.7cm]
&&\rnode{b}{$-1$}&&&\\[0.7cm]
\ncline[linewidth=0.7pt,linecolor=yellow]{-}{a}{b}\bput[1pt]{0}(0.5){\small$J$}%
\ncline[linewidth=0.7pt,linecolor=yellow]{-}{a}{c}\aput[1pt]{0}(0.5){\small$S$}%
%&\circlenode*{b}{\makebox[0.5cm]{{\red$B^{\blue11}$}}}&\circlenode*{e}{\makebox[0.5cm]{{\red$E^{\blue4}$}}}&&\\[0.1cm]
%&&&\circlenode*{h}{\makebox[0.5cm]{{\red$H^{\blue3}$}}}&\\[0.1cm]
%\circlenode*{a}{\makebox[0.5cm]{{\red$A^{\blue11}$}}}&\circlenode*{c}{\makebox[0.5cm]{{\red$C^{\blue7}$}}}%
\end{tabular}

On voit donc qu'il faut choisir $S_1$ au premier n{\oe}ud. La
procédure d'induction à rebours aboutit donc à la stratégie
$S_1J_2$. C'est parce qu'on a conclut qu'il fallait jouer $J_2$ et
que le gain correspondant est de $-1$ qu'on a pu ensuite conclure
qu'il fallait jouer $S_1$.

\section{Décision dynamique en information imparfaite}

On suppose maintenant que l'agent n'est pas complètement informé
de tous les coups de la nature. Au moment où il effectue ses
choix, il n'a qu'une information imparfaite sur sa position dans
l'arbre de décision. Nous commençons par décrire son information
dans l'arbre à l'aide de la notion d'ensembles d'information.

Un ensemble d'information est une classe d'équivalence sur les
n{\oe}uds $t$ tels que $i(t)=J$. Soit $h(t)$ la classe
d'équivalence du n{\oe}ud $t$, c'est-à-dire l'ensemble
d'information comprenant le n{\oe}ud $t$. Comme $h$ décrit des
relations d'équivalence on a~:
$$
t' \in h(t) \implies t \in h(t')
$$
 On suppose que l'agent conna\^it son ensemble de décisions à chaque n{\oe}ud.
L'ensemble d'information $h(t)$. On a la restriction suivante sur
la fonction qui à $t$ associe $h(t)$ :
$$
t'\in h(t) \implies A(t')=A(t)
$$

\begin{exemple}
Reprenons l'exemple du casino en supposant qu'après la première
étape, le joueur ne sait pas s'il a gagné ou non. On a l'arbre de
décision suivant :

\begin{tabular}{p{1cm}p{1cm}p{1cm}p{1cm}p{1cm}p{1cm}}
&&&\rnode{a}{$\bullet$}&\rnode{c}{$0$}\\[0.7cm]
&&\rnode{b}{$\blue\diamond$}&&&\\[0.7cm]
&\rnode{d}{$\bullet$}&&\rnode{e}{$\bullet$}&\\[0.7cm]
&\rnode{f}{$\blue\diamond$}&\rnode{g}{\small$+1$}&\rnode{h}{$\blue\diamond$}&\rnode{i}{\small$-2$}&\\[0.7cm]
\rnode{j}{\small$+4$}&\rnode{k}{\small$-1$}&\rnode{l}{\small$+1$}&\rnode{m}{\small$-4$}\\[0.7cm]
\ncline[linewidth=0.7pt,linecolor=yellow]{-}{a}{b}\bput[1pt]{0}(0.5){\small$J$}%
\ncline[linewidth=0.7pt,linecolor=yellow]{-}{a}{c}\aput[1pt]{0}(0.5){\small$S$}%
\ncline[linewidth=0.7pt,linecolor=yellow]{-}{b}{d}\bput[1pt]{0}(0.5){\small$\frac12$}%
\ncline[linewidth=0.7pt,linecolor=yellow]{-}{b}{e}\aput[1pt]{0}(0.5){\small$\frac12$}%
\ncline[linewidth=0.7pt,linecolor=yellow]{-}{d}{f}\bput[1pt]{0}(0.5){\small$J$}%
\ncline[linewidth=0.7pt,linecolor=yellow]{-}{d}{g}\aput[1pt]{0}(0.5){\small$S$}%
\ncline[linewidth=0.7pt,linecolor=yellow]{-}{e}{h}\bput[1pt]{0}(0.5){\small$J$}%
\ncline[linewidth=0.7pt,linecolor=yellow]{-}{e}{i}\aput[1pt]{0}(0.5){\small$S$}%
\ncline[linewidth=0.7pt,linecolor=yellow]{-}{f}{j}\bput[1pt]{0}(0.5){\small$\frac12$}%
\ncline[linewidth=0.7pt,linecolor=yellow]{-}{f}{k}\aput[1pt]{0}(0.5){\small$\frac12$}%
\ncline[linewidth=0.7pt,linecolor=yellow]{-}{h}{l}\bput[1pt]{0}(0.5){\small$\frac12$}%
\ncline[linewidth=0.7pt,linecolor=yellow]{-}{h}{m}\aput[1pt]{0}(0.5){\small$\frac12$}%
\psellipse[linecolor=green](2.9,3.66)(1.8,0.3)
%&\circlenode*{b}{\makebox[0.5cm]{{\red$B^{\blue11}$}}}&\circlenode*{e}{\makebox[0.5cm]{{\red$E^{\blue4}$}}}&&\\[0.1cm]
%&&&\circlenode*{h}{\makebox[0.5cm]{{\red$H^{\blue3}$}}}&\\[0.1cm]
%\circlenode*{a}{\makebox[0.5cm]{{\red$A^{\blue11}$}}}&\circlenode*{c}{\makebox[0.5cm]{{\red$C^{\blue7}$}}}%
\end{tabular}

Un premier ensemble d'information est donné par le n{\oe}ud au
début du jeu. Un deuxième contient les deux n{\oe}uds qui suivent
le choix de la nature après une première décision de jouer.
\end{exemple}

Une stratégie $s$ est une fonction de l'ensemble des {\em
ensembles d'information} vers les décisions possibles. On peut
voir $s$ comme une application de $i^{-1}(J)$ vers $\cup_t A(t)$
telle que :%
\begin{itemize}
\item $s(t)= s(t')$ si $h(t) = h(t')$ ($s(t)$ ne dépend que de
l'ensemble d'information de $t$)%
\item $s(t)\in A(t)$ pour tout $t$.
\end{itemize}

\begin{exemple}
On a deux ensembles d'information, et deux décisions possibles
pour chacun d'entre eux, soient 4 stratégies. $JS$, $JJ$, $SJ$ et
$SS$. Les paiement espérés correspondants sont :
\begin{itemize}
\item $JS$ donne $-\frac12$%
\item $JJ$ donne $0$%
\item $SJ$ et $SS$ donnent $0$.
\end{itemize}
\end{exemple}

\begin{exemple}
Rechercher les stratégies qui maximisent la probabilité de gagner
au moins 1 dans ce dernier jeu.
\end{exemple}

\chapter{Jeux sous forme normale}
La théorie de la décision --par exemple du consommateur-- permet
de relier les choix d'un agent avec ses préférences. Un agent
rationnel est un agent qui cherche à maximiser son utilité
(espérée). Un ``bon'' choix, ou choix rationnel, est un choix
maximisateur d'utilité.

Nous allons introduire la théorie des jeux, que l'on pourrait
aussi appeler théorie de la décision interactive, en modélisant et
en étudiant des situations dans lesquelles plusieurs agents sont en
situations de choix, les choix des uns affectant l'utilité des
autres. Le ``bon'' choix pour un agent dépend maintenant des choix
des autres agents.

\section{Description d'un jeu sous forme normale}
\begin{exemple}
Amanda et Bertrand veulent se rencontrer à NY, mais n'ont pas
décidé d'un lieu de rendez-vous. Ils peuvent chacun se rendre soit
à l'Empire State Building, soit à Central Park.

Le but de chacun est de rencontrer l'autre, peu importe où. Les
préférences de chacun sur quel lieu o\`u aller dépendent donc de
ce que fait l'autre.

On peut donc décrire la situation par les éléments suivants :
\begin{itemize}
\item 2 joueurs, Amanda et Bertrand %
\item 2 décisions possibles pour chacun d'entre eux. (ESB, CP).%
\item L'utilité de chaque joueur en fonction de son choix et de celui
de l'autre joueur.
\end{itemize}
\end{exemple}

En général, un jeu sous forme normale est donné par
\begin{itemize}
\item  Une liste d'agents (ou joueurs) $I$; %
\item  Un ensemble de stratégies $S_i$ pour chaque agent $i$;%
\item  L'issue du jeu $s$ correspond au vecteur $(s_i)$ de
stratégies choisi; %
\item Des préférences pour les joueurs sur $S$,
représentées par une fonction $g_i \colon \Pi_i S_i\to \reals$.
\end{itemize}

\begin{exemple} Pour le jeu de rendez-vous à New-York,
\begin{itemize}
\item  $I=\{\mbox{Amanda},\mbox{ Bernard}\}$; %
\item  $S_i = \{E, C\}$; %
\item  $S = \{(E,C), (E,E), (C,C),
(C,E)\}$; %
\item $g_i(s_1, s_2) = 1$ si $s_1 = s_2$, $0$ sinon.
\end{itemize}
On représente le jeu par la matrice suivante : Amanda choisit la
ligne, et Bernard choisit la colonne. La paire de paiements pour
Amanda et pour Bernard est inscrite dans la case.


\begin{center}
\begin{game}{2}{2}
    &  $E$  &  $C$    \\
$E$ & $1,1$ & $0,0$   \\
$C$ & $0,0$ & $1,1$   \\
\end{game}\hspace*{20mm}%
\end{center}

\end{exemple}

\begin{definition}
Un jeu $G$ sous forme normale est donné par :
\begin{itemize}
\item Un ensemble fini de joueurs $I$ ; %
\item Des ensembles de
stratégies $S_i$ ; %
\item Des fonctions de paiement $g_i\colon
\Pi_i S_i\to \reals$.
\end{itemize}
On utilise les notations :\\
$S=\Pi_i S_i$ \\
$S_{-i} = \Pi_{j\neq i}S_j$.
\end{definition}

Nous passons maintenant en revue quelques jeux ou classes de jeux
particulièrement importants.

\subsection{Jeux à intérêts communs}
Ce sont des jeux dans lesquels tous les joueurs ont les mêmes
intérêts, ou préférences. On a donc $g_i = g_j$ pour tous joueurs
$i,j$. On peut voir ces jeux comme une généralisation des
problèmes à un agent.

Par exemple, le jeu de rendez-vous à New-York fait partie de cette
catégorie.

\subsection{Jeux à somme nulle}
Ce sont des jeux à deux joueurs dans lesquels les intérêts sont
parfaitement antagonistes. Ces jeux sont à l'opposé des jeux à
intérêts communs. Ce qui est gagné par un joueur est perdu par
l'autre. La somme des fonctions d'utilité est donc $0$. C'est à
dire que $g_i = -g_j$.

\begin{exemple} Jeux de pile ou face. Deux joueurs annoncent
Pile (P) ou face (F). Le joueurs 2 paie 1 \euro\ à 1 si les deux
choisissent la même chose, sinon le joueur 1 paie 1 \euro\ au
joueur 2.\\
\begin{center}
\begin{game}{2}{2}
    &  $P$  &  $F$    \\
$P$ & $1,-1$ & $-1,1$   \\
$F$ & $-1,1$ & $1,-1$   \\
\end{game}\hspace*{20mm}%
\end{center}
Chaque case représente l'utilité du joueur 1, suivie par celle du
joueur~2.
\end{exemple}

\subsection{La bataille des sexes}
Certains jeux font intervenir une part de coordination et une part de conflit entre
les agents. C'est le cas du jeu de la bataille des sexes suivant.\\


\begin{exemple}
Un couple veut décider d'une sortie.\\
L'homme préfère aller à l'Opéra, et
la femme préfère assister à un match de foot.\\
Pour chacun, être avec l'autre est plus important que le lieu.\\
\begin{center}
\begin{game}{2}{2}
    &  $F$  &  $O$    \\
$F$ & $2,1$ & $0,0$   \\
$O$ & $0,0$ & $1,2$   \\
\end{game}\hspace*{20mm}%
\end{center}
\end{exemple}

\subsection{La fureur de vivre} Deux adolescents en voiture foncent
l'un vers l'autre sur une route étroite. Personne ne veut sortir
de la route. Si les deux sortent, aucun n'est vraiment satisfait,
ni mécontent. Les choix pour chacun sont $F$ (faucon, rester sur
la route), ou $C$ (colombe, sortir de la route).
\begin{center}
\begin{game}{2}{2}
    &  $F$  &  $C$    \\
$F$ & $-1,-1$ & $10,0$   \\
$C$ & $0,10$ & $5,5$   \\
\end{game}\hspace*{20mm}%
\end{center}

\subsection{Dilemme du prisonnier} Deux prisonniers complices sont
interrogés séparément. Chacun peut trahir son partenaire en le
dénonçant ($T$), ou bien rester silencieux ($S$). Si les deux
trahissent, ils vont en prison pour 5 ans chacun. Si l'un trahit
et pas l'autre celui qui trahit sort libre et l'autre va en prison
pour 10 ans. Si personne ne trahit, ils vont en prison pour 3 ans
tous deux.
\begin{center}
\begin{game}{2}{2}
    &  $S$  &  $T$    \\
$S$ & $-3,-3$ & $-10,0$   \\
$T$ & $0,-10$ & $-5,-5$   \\
\end{game}\hspace*{20mm}%
\end{center}

Le jeu du dilemme du prisonnier est un exemple fondamental en
économie. Ce jeu est important car il fait ressortir une tension
entre l'intérêt individuel et l'intérêt collectif. De nombreuses
situations présentent une structure similaire à celle du dilemme
du prisonnier
\begin{itemize}
\item Achat par internet ; Un acheteur et un vendeur se sont mis
d'accord sur internet. Chacun a le choix entre envoyer le colis
(ou bien l'argent), et ne pas l'envoyer. Chacun préfère que
l'autre l'envoie et préférerait ne pas envoyer sa part. Cependant,
les deux sont plus satisfaits si chacun respecte sa
part du contrat plutôt que si aucun ne le fait.%
\item Course aux armements ; Chaque pays peut décider de s'armer,
ou non. L'intérêt des deux pays est qu'aucun ne dépense de
resources pour s'armer, mais à stratégie fixée de l'autre, chacun
préfère s'armer.%
\item Achat d'un 4x4 ; Un 4x4 est avantageux pour celui qui l'a car on peut impressionner
les autres voitures, et on se sent plus en sécurité. Mais ceci est au détriment des autres voitures. %
\item Collusion et la commission européenne ; La commission
européenne cherche à lutter contre les ententes secrètes des
industries. Le problème étant que personne ne veut dénoncer ces
ententes. La règle est que celui qui dénonce une entente ne se
verra pas poursuivi. Chacun a donc intérêt à dénoncer plutôt qu'à
être dénoncé.
\end{itemize}


\subsection{Compétition en quantités dite de Cournot}
Deux firmes, 1 et 2, produisent des biens identiques.\\
Chacune décide d'un niveau de production $q_i$, à un co\^ut
$c_i(q_i)$. Le prix résultant de la loi de l'offre et de la
demande est $p(q_1+q_2)$.
\begin{eqnarray*}
I   &=& \{1,2\}\\
S_i &=& \reals_+ \\
g_1(q_1,q_2) &=& q_1 p(q_1+q_2)-c_1(q_1)\\
g_2(q_1,q_2) &=& q_2 p(q_1+q_2)-c_2(q_2)
\end{eqnarray*}

\subsection{Compétition en prix dite de Bertrand}
Il s'agit d'un modèle de compétition par les prix. Chaque firme
décide d'un prix de vente du bien, et les consommateurs (une unité
en tout) achètent à la firme au prix le plus bas.
\begin{eqnarray*}
I   &=& \{1,2\}\\
S_i &=& [0,M] \\
\end{eqnarray*}
$$
g_i(p_i,p_j)=\begin{cases}
  p_i &\mbox{si }\;\; p_i < p_j\\
  0 &\mbox{si }\;\; p_i > p_j\\
  p_i/2 &\mbox{si }\;\; p_i = p_j
\end{cases}$$\\

\section{Stratégies dominantes et dominées}
Nous commençons maintenant l'étude des prédictions des issues des
jeux sous forme normale. Étant donné un jeu, et en supposant les
joueurs rationnels, que peut-on prédire qu'ils vont jouer, ou que
peut-on prédire qu'ils ne vont pas jouer ? Si nous devions jouer
dans un tel jeu, quel choix ferions-nous ?

Une stratégie dominée est une stratégie qui donne un paiement
strictement moins bon que celui d'une autre stratégie donnée, ceci
pour tout choix possible des adversaires.

\begin{definition}
Une stratégie $s_i$ est (strictement) dominée s'il existe $s'_i$
telle que
$$
\forall s_{-i}\in S_{-i},\, g_i(s'_i,s_{-i}) > g_i(s_i,s_{-i})
$$
On dit alors que $s'_i$ domine (strictement) $s_i$.
\end{definition}

Un joueur ``rationnel'' ne devrait jamais jouer une stratégie
dominée. En effet, il a la possibilité d'une autre stratégie dont
il sait qu'elle peut lui donner un meilleur gain, indépendamment
des choix des autres joueurs.

Une stratégie dominante donne un meilleur paiement que toute autre
stratégie, pour tous les choix des autres joueurs.

\begin{definition}
Une stratégie $s_i$ est (strictement) dominante si pour toute
autre stratégie $s'_i$.
$$
\forall s_{-i}\in S_{-i},\, g_i(s_i,s_{-i}) > g_i(s'_i,s_{-i})
$$
\end{definition}

La définition se réécrit :
\begin{proposition} Une stratégie $s_i$ est dominante si et seulement si elle domine toute stratégie
$s'_i\neq s_i$.
\end{proposition}

Une conséquence quasi-immédiate de la définition est : %
\begin{proposition}
Si une stratégie $s_i$ est dominante, alors elle est unique à
avoir cette propriété.
\end{proposition}

Si un joueur a une stratégie dominante, on peut alors penser que
cette stratégie constitue un bon choix. Pour chaque choix des
autres, elle donne le meilleur paiement possible.

Nous pouvons maintenant construire un raisonnement en se basant
sur le fait que si aucun joueur ne choisit de stratégie dominée,
alors les autres joueurs peuvent anticiper ce phénomène.

\begin{exemple}
Considérons le jeu suivant :
\begin{center}
\begin{game}{2}{2}
    &  $S$  &  $T$    \\
$S$ & $0,-2$ & $-10,-1$   \\
$T$ & $-1,-10$ & $-5,-5$   \\
\end{game}\hspace*{20mm}%
\end{center}

Que devrait jouer le joueur 1 s'il sait que le joueur 2 est
rationnel ?
\end{exemple}

\begin{exemple}
Considérons le jeu suivant :
\begin{center}
\begin{game}{2}{3}
    &  $G$  &  $M$   & $D$    \\
$H$ & $2,2$ & $1,1$  & $4,0$  \\
$B$ & $1,2$ & $4,1$  & $3,5$  \\
\end{game}\hspace*{20mm}%
\end{center}

\begin{itemize}
\item Que devrait jouer le joueur 1 s'il sait que le joueur 2 est rationnel ?\\
\item Que devrait jouer le joueur 2 s'il sait que le joueur 1 sait qu'il est rationnel ?\\
\end{itemize}

\end{exemple}

Ceci nous amène à définir le processus d'élimination itérée des stratégies dominées.

\begin{definition} La procédure d'élimination itérée des
stratégies dominées est la suivante :
\begin{itemize}
\item {\bf \'Etape 1 :} $S_i^0 = S_i$.%
\item {\bf \'Etape 2 :} $S_i^1 = \{$ stratégies non dominées
de $i$ $\}$.%
\item {\bf \'Etape $k+1$ :} $$S_i^{k+1} = \{s_i \in S_i^k,
\not\exists s'_i \in S^{k}_i, \forall s_{-i}\in
S^{k}_{-i}, g_i(s'_i,s_{-i})>g_i(s_i,s_{-i})\}$$%
Stratégies non dominées {\em face aux stratégies de } $S^k_i$%
\item {\bf \'Etape $\infty$ :} $S_i^{\infty} = \cap_k S^k_i$.
\end{itemize}
\end{definition}

Pour certains jeux, cette définition nous conduit à une prédiction
unique des stratégies suivies par les joueurs.

\begin{definition}
Un jeu $G$ est résoluble par élimination itérée de stratégies
dominées si pour tout $i$, $S_i^{\infty}$ est un singleton.
\end{definition}

\begin{exemple}
Que dit la procédure d'élimination de stratégies dominées dans le
jeu suivant ?
\begin{center}
\begin{game}{4}{4}
    &  $A$  &  $B$   & $C$   & $D$  \\
$A$ & $5,2$ & $2,6$  & $1,4$ & $0,4$\\
$B$ & $0,0$ & $3,2$  & $2,1$ & $1,1$\\
$C$ & $7,0$ & $2,2$  & $1,5$ & $5,1$\\
$D$ & $9,5$ & $1,3$  & $0,2$ & $4,8$\\
\end{game}\hspace*{20mm}%
\end{center}
\end{exemple}


Passons maintenant à l'exemple suivant.
\begin{exemple}~\\
\begin{center}
\begin{game}{2}{2}
    &  $G$  &  $D$\\
$H$ & $2,2$ & $4,4$\\
$B$ & $3,1$ & $4,1$\\
\end{game}\hspace*{20mm}%
\end{center}
Quelles sont les stratégies dominées, dominantes ?\\
Quelles stratégies paraissent ``raisonnables'' ?%
\end{exemple}

Dans l'exemple précédent, il n'y a pas de stratégies (strictement)
dominantes, ni dominées. Cependant, le choix de la stratégie $B$
pour le joueur 1, et $D$ pour le joueur 2, semble un bon choix,
car ces stratégies ne peuvent donner qu'un meilleur paiement que
l'autre, et jamais un paiement moins bon. Ceci nous conduit à la
définition d'une stratégie faiblement dominée.

\begin{definition} Une stratégie $s_i$ est faiblement
dominée s'il existe $s'_i$ telle que:
\begin{eqnarray*}
\forall s_{-i}, g_i(s'_i,s_{-i}) &\geq& g_i(s_i,s_{-i})\\
\exists s_{-i}, g_i(s'_i,s_{-i}) &>& g_i(s_i,s_{-i})
\end{eqnarray*}
On dit dans ce cas que $s'_i$ domine faiblement $s_i$.
\end{definition}

Dans le jeu précédent, $H$ est faiblement dominée pour le joueur
1, et $G$ est faiblement dominée pour le joueur 2.

\begin{definition}
Une stratégie $s_i$ est faiblement dominante si pour toute autre
stratégie $s'_i$.
\begin{eqnarray*}
\forall s_{-i}\in S_{-i},\, g_i(s_i,s_{-i}) &\geq &g_i(s'_i,s_{-i})\\
\exists s_{-i}\in S_{-i},\, g_i(s_i,s_{-i}) &> &g_i(s'_i,s_{-i})
\end{eqnarray*}
\end{definition}

Si une telle stratégie existe, elle est unique. Elle semble
représenter un choix ``raisonnable'', car on ne peut jamais le
regretter.

La définition se réécrit :
\begin{proposition} Une stratégie $s_i$ est faiblement dominante si et seulement
si elle domine faiblement toute stratégie $s'_i\neq s_i$.
\end{proposition}

Une conséquence quasi-immédiate de la définition est : %
\begin{proposition}
Si une stratégie $s_i$ est faiblement dominante, alors elle est
unique à avoir cette propriété.
\end{proposition}

\begin{exemple}
Un groupe de joueurs $I$ participe à une enchère pour un objet. La
valeur de cet objet est $v_i$ pour le joueur $i$. Si chaque joueur
$i$ met une enchère de $e_i$, l'objet est attribué à celui qui met
l'enchère la plus haute, à un prix correspondant à la seconde
enchère la plus haute.

On demande de montrer comme exercice que pour le joueur $i$, jouer
$e_i = v_i$ est une stratégie faiblement dominante.
\end{exemple}

Lorsqu'un joueur possède une stratégie dominante, on peut penser
que cette stratégie constitue un ``bon'' choix. En revanche, une
stratégie dominée (même faiblement) ne semble pas constituer un
bon choix. Nous sommes tentés de généraliser la procédure de
suppression itérée de stratégies dominées aux stratégies
faiblement dominées. Voyons ce que l'on obtient dans l'exemple
suivant :

\begin{exemple} Soit le jeux donné par la matrice de paiements

\begin{center}
\begin{game}{3}{2}
    &  $G$  &  $D$\\
$H$ & $1,1$ & $0,0$\\
$M$ & $1,1$ & $2,1$\\
$B$ & $0,0$ & $2,1$\\
\end{game}\hspace*{20mm}%
\end{center}

La procédure d'élimination des stratégies faiblement dominées nous
permet d'éliminer $T$, puis $L$. Cette procédure pourrait aussi
bien nous conduire à éliminer $B$, puis $R$.  Ces deux ordres de
suppression de stratégies dominées nous conduit à des résultats
différents. La méthode n'est donc pas concluante.
\end{exemple}

On demande de montrer en exercice que l'ordre de suppression des
stratégies strictement dominées n'a pas d'influence sur le
résultat final.

En conclusion : Si on supprime de manière itérée les stratégies
{\bf faiblement} dominées, le résultat peut dépendre de l'ordre
des itérations. Ce n'est donc pas une bonne procédure. En
revanche, nous n'avons pas ces difficultés avec la suppression
itérée de stratégies strictement dominées.

\begin{exemple}
Considérons le jeu de compétition de Cournot suivant :
\begin{itemize}
\item $c_i(q_i) = c q_i$. \item $p = \alpha - \beta(q_1+q_2)$.
\end{itemize}
Si $i$ croit que $j$ va choisir $q_j$, alors la meilleure réponse
(la stratégie qui maximise le paiement de $i$ à stratégie de $j$
fixée) est :
$$
MR_i(q_j)=
\begin{cases}
  \frac{\alpha-c}{2\beta} - \frac{q_j}2 &\mbox{si }\;\; q_i \leq\frac{\alpha-c}{\beta}\\
   0 &\mbox{sinon }
\end{cases}$$\\
Les joueurs choisissent dans $S^0=\reals$.\\
Toutes les stratégies hors de $S^1=[0,\frac{\alpha-c}{2\beta}]$
sont dominées
(par 0).\\
Ensuite, toutes les stratégies hors de
$S^2=[MR(\frac{\alpha-c}{2\beta}), MR(0)]$ sont dominées (par
$MR(\frac{\alpha-c}{2\beta})$). En itérant, on démontre que le jeu
est résoluble par élimination itérée de stratégies dominées et
que, $S_i^{\infty}=\{\frac{\alpha-c}{3\beta}\}$.
\end{exemple}

\section{\'Equilibre de Nash}
La suppression itérée de stratégies dominées est attractive car
elle repose sur l'idée de rationalité (et sur la connaissance par
les joueurs de la rationalité des autres, etc...). Cependant, peu
de jeux sont résolubles par élimination itérée de stratégies
dominées. Nous introduisons un concept qui s'applique à une plus
grande gamme de jeux : l'équilibre de Nash.

L'équilibre de Nash peut être compris comme une convention sociale
stable. C'est une règle collective de laquelle aucun individu n'a
intérêt à s'écarter pourvu que les autres individus respectent eux
aussi la règle.

\begin{definition} Un profil de stratégies $s^*=(s_i)_i$ est
un équilibre de Nash (en stratégies pures) si :
$$
\forall i, \forall s'_i\in S_i, g_i(s_i,s_{-i})\geq
g_i(s'_i,s_{-i})
$$
\end{definition}

La notion suivante d'équilibre de Nash strict est plus forte. Elle
suppose que les joueurs préfèrent tous strictement jouer
l'équilibre de Nash plutôt qu'une autre stratégie si les autres
joueurs suivent l'équilibre de Nash.
\begin{definition}  Un profil de stratégies $s^*\in S$ est un
équilibre de Nash {\em strict} (en stratégies pures) si
$$
\forall i, \forall s'_i \in S_i, g_i(s^*_i,s^*_{-i})>
g_i(s'_i,s^*_{-i})
$$
\end{definition}


En particulier, un équilibre de Nash strict est un équilibre de
Nash.

\begin{exemple} Dilemme du prisonnier
\begin{center}
\begin{game}{2}{2}
    &  $S$  &  $T$\\
$S$ & $-3,-3$ & $-10,-10$\\
$T$ & $0,-10$ & $-5,-5$\\
\end{game}\hspace*{20mm}%
\end{center}
Seul équilibre de Nash: $(T,T)$.
\end{exemple}

\begin{exemple} Soit le jeu
\begin{center}
\begin{game}{2}{3}
    &  $G$  &  $M$  &  $D$\\
$H$ & $2,2$ & $1,1$ & $4,0$\\
$B$ & $1,2$ & $4,1$ & $3,5$\\
\end{game}\hspace*{20mm}%
\end{center}
Unique équilibre de Nash : $(H,G)$.
\end{exemple}

Le jeu du dilemme du prisonnier est résoluble en stratégies
dominées, et l'équilibre de Nash est compatible avec cette
prédiction.

Le jeu précédent n'est pas résoluble en stratégies dominées, et
cependant l'équilibre de Nash nous donne une prédiction.

A l'équilibre de Nash, aucune stratégie dominée n'est jouée. Plus
généralement, toutes les stratégies des équilibres de Nash
survivent à l'élimination itérée des stratégies dominées.

\begin{proposition}
Si $s^*$ est un équilibre de Nash, alors pour tout $i$, $s^*_i\in
S^\infty_i$.
\end{proposition}
{\em Démonstration :} Supposons $s\in \Pi_i S^k_i$, $s_i\not\in
S^{k+1}_i$. Alors il existe $s'_i$ meilleur pour $i$ que $s_i$
contre tous $s_{-i}\in \Pi_{j\neq i} S^k_j$, donc en particulier
contre $s^*_{-i}$. Contradiction.

Si on a un profil de stratégies dominantes, ce profil constitue un
équilibre de Nash:
\begin{proposition}
Si $s^*$ est un profil de stratégies faiblement dominantes, $s^*$
est un équilibre de Nash. Si $s^*$ est un profil de strictement
dominantes alors $s^*$ est l'unique
équilibre de Nash. %
\end{proposition}

\begin{exemple} Le jeu suivant est-il résoluble
par élimination itérée de stratégies dominées ?\\
Quels sont ses équilibres de Nash ?

\begin{center}
\begin{game}{3}{3}
    &  $G$  &  $M$  &  $D$ \\
$h$ & $5,3$ & $0,4$ & $3,5$\\
$m$ & $4,0$ & $5,5$ & $4,0$\\
$b$ & $3,5$ & $0,4$ & $5,3$\\
\end{game}\hspace*{20mm}%
\end{center}
\end{exemple}

\begin{exemple}
Reprenons le jeu de compétition de Cournot donné par :
\begin{itemize}
\item $c_i(q_i) = c q_i$. \item $p = \alpha - \beta(q_1+q_2)$.
\end{itemize}
Quels sont les équilibres de Nash ?
\end{exemple}

\begin{exemple}
Deux firmes, 1, 2. Firme $i$ décide d'un prix $p_i$. Le co\^ut
marginal de production est $c$. La demande est donnée par $D(p)$,
tous les acheteurs achêtent au meilleur prix.
$$
D_i(p_1,p_2)=
\begin{cases}
  D(p_i)   &\mbox{si }\;\; p_i < p_j\\
  D(p_i)/2 &\mbox{si }\;\; p_i = p_j\\
  0        &\mbox{si }\;\; p_i > p_j
\end{cases}$$\\
On suppose $D$ décroissante, et $D(c)>0$.\\

Quels sont les équilibres de Nash ?
\end{exemple}

\begin{exemple}
Modèle de compétition spatiale, appelée aussi compétition de Hotelling.\\
Deux vendeurs de glace sur une plage linéaire ([0,1]).\\
Chacun choisit un emplacement $x_i\in [0,1]$. Les consommateurs
sont uniforméments répartis sur la plage,
et achètent au vendeur le plus proche.\\
Chaque vendeur veut maximiser ses ventes.\\
Représentation en jeux ?\\
\'Equilibres de Nash ?

\end{exemple}

Un jeu peut avoir plusieurs équilibres de
Nash.

\begin{exemple} Quels sont les équilibres de Nash en stratégies pures
du jeu du RdV à NY ?\\
\begin{center}
\begin{game}{2}{2}
    &  $E$  &  $C$  \\
$E$ & $1,1$ & $0,0$ \\
$C$ & $0,0$ & $1,1$ \\
\end{game}\hspace*{20mm}%
\end{center}
\end{exemple}

On peut aussi ne pas avoir d'équilibres. Voir le jeu suivant.
\begin{exemple}
Quels sont les équilibres de Nash en stratégies pures
du jeu suivant à somme nulle ?\\
\begin{center}
\begin{game}{2}{2}
    &  $P$  &  $F$  \\
$P$ & $-1,1$ & $1,-1$\\
$F$ & $1,-1$ & $-1,1$\\
\end{game}\hspace*{20mm}%
\end{center}
\end{exemple}

On peut avoir plusieurs équilibres de Nash, et cependant en
préférer un plutôt qu'un autre pour des raisons liées au
``risque'' pris en jouant les stratégies.

\begin{exemple}
Quels sont les équilibres de Nash en stratégies pures du jeu
suivant ?\\ Quelle stratégie choisiriez-vous ?
\begin{center}
\begin{game}{2}{2}
    &  $A$  &  $B$  \\
$A$ & $9,9$ & $-15,8$\\
$B$ & $8,-15$ & $7,7$\\
\end{game}\hspace*{20mm}%
\end{center}
\end{exemple}

Certaines stratégies faiblement dominées peuvent être jouées à
l'équilibre de Nash.

\begin{exemple}
Quels sont les équilibres de Nash en stratégies pures du jeu
suivant ?\\ Quelle stratégie choisiriez-vous ?
\begin{center}
\begin{game}{2}{2}
    &  $A$  &  $B$  \\
$A$ & $2,2$ & $1,2$\\
$B$ & $2,1$ & $1,1$\\
\end{game}\hspace*{20mm}%
\end{center}
\end{exemple}

\subsection{Interprétations de l'\'Equilibre de Nash}
La notion d'équilibre de Nash est centrale en théorie des jeux et
en microéconomie. Elle est aussi appliquée en informatique,
biologie, en sociologie... Les interprétations de cette notion
sont diverses. On peut citer les justifications suivantes :
\begin{itemize} %
\item Prescriptions. Un individu extérieur
propose un profil de stratégies. Si ce profil est un équilibre
de Nash, personne n'a intérêt à dévier.%
\item Communication. Les joueurs se mettent d'accord à l'avance
sur quel profil jouer. Les EN sont des accords viables.%
\item Introspection. Un joueur analyse le jeu et se convainc que
seul l'EN est possible. Valable uniquement si unique EN. %
\item Norme sociale. Dans certains pays, nous conduisons à
droite de la route, dans d'autres à gauche. C'est la norme sociale
qui dicte une règle. EN est nécessaire pour que cette
règle soit suivie.%
\item Apprentissage. Les joueurs découvrent les stratégies
suivies par les autres au cours d'un processus d'apprentissage.%
\item Evolution.  Si chaque espèce évolue de fa\c{c}on
adaptée aux autres espèces, la résultante est un EN.
\end{itemize}

\chapter{Stratégies mixtes}
Un des problèmes inhérents au concept d'équilibre de Nash en
stratégies pures est que pour certains jeux, de tels équilibres
n'existent pas. P.ex.\,le jeu de ``Pierre, Papier, Ciseaux" :
\begin{center}
\begin{game}[Pierre, Papier, Ciseaux]{3}{3}
     &  $Pi$    &  $Pa$  & $Ci$\\
$Pi$ & $0,0$    & $-1,1$ & $1,-1$\\
$Pa$ & $1,-1$   & $0,0$  & $-1,1$\\
$Ci$ & $-1,1$   & $1,-1$ & $0,0$\\
\end{game}\hspace*{20mm}%
\end{center}
n'admet pas d'équilibre de Nash.

La raison pour laquelle on n'a pas d'équilibres est la suivante :
la notion d'équilibre de Nash en stratégies pures suppose que
chaque joueur connaît les stratégies des autres joueurs. Or, nous
sommes dans des jeux ou chaque joueur a intérêt à cacher sa
stratégie, ou a bluffer. En effet, dans les jeux ``pile ou face"
ou ``tirer un penalty", ou ``bluff au poker", on n'utilise pas
toujours la même stratégie, et on ne connaît pas non plus à
l'avance celle de l'adversaire.
\begin{center}
\begin{game}[Pile ou face]{2}{2}
     &  $Pi$  &  $Fa$\\
$Pi$ & $1,-1$ & $-1,1$\\
$Fa$ & $-1,1$ & $1,-1$
\end{game}\hspace*{20mm}%
\end{center}

Les stratégies mixtes, ou aléatoires, vont permettre de
représenter ces possibilités de bluff, ou de jouer aléatoirement.

\section{Définition des concepts}
Nous allons définir les stratégies mixtes, puis les jeux dans
lesquels les joueurs ont la possibilité de jouer des stratégies
mixtes, et enfin définir un équilibre de Nash en stratégies mixtes
comme un équilibre de Nash d'un tel jeu.

\subsection{Stratégies mixtes}
Nous commençons par définir le concept de stratégies mixtes. Soit
donc $G = (I, (S_i)_i, (g_i)_i)$ un jeu sous forme normale

\begin{definition} Une strat\'egie mixte pour le joueur
$i$ est une loi de probabilit\'es sur $S_i$. $\Sigma_i
=\Delta(S_i)$ repr\'esente l'ensemble des strat\'egies
mixtes du joueur $i$. %
\end{definition}

On utilisera les notations : $\Sigma = \Pi_i \Sigma_i$,
$\Sigma_{-i} = \Pi_{j\neq i} \Sigma_j$.

Par opposition aux stratégies mixtes, on appelle les éléments de
$S_i$ {\bf stratégies pures}. Ce sont des stratégies
déterministes.

On utilisera selon les cas différentes notations pour une
stratégie mixte.
\begin{enumerate}
\item Notation fonction. $\sigma_i$ est l'application de $S_i$
vers $\reals$ qui associe à la stratégie pure $s_i$ sa probabilité
d'être jouée. Exemple: $S_1=\{H,B\}$,
$\sigma_1(H)=\sigma_1(B)=\frac12$%
\item Vecteur. On écrit $\sigma_i$ comme un vecteur de
probabilités. Si $S_i= \{s_{i,1},\ldots,s_{i,n_i}\}$, on
$\sigma_i$ se représente comme
$(\sigma_i(s_{i,1}),\ldots,\sigma_i(s_{i,n_i}))$. Exemple :
$\sigma_1=(\frac12,\frac12)$. %
\item Combinaison convexe de stratégies pures : $\frac12H +
\frac12B$.
\end{enumerate}

Il y a plusieurs interprétations possibles des stratégies mixtes,
parmi elles, les plus courantes sont les suivantes :
\begin{itemize}
\item Possibilit\'e ``mat\'erielle'' de jouer al\'eatoirement.
Rien ne s'oppose à priori à ce qu'un joueur décide aléatoirement
quelle stratégie pure utiliser. La stratégie mixte est un choix
aléatoire d'une stratégie pure.%
\item Bluff, et part d'incertitude que chacun laisse sur sa
strat\'egie. Ici, une stratégie mixte représente plutôt une
croyance que chaque joueur a sur les manières de jouer des autres
joueurs. Même si je décide si je vais bluffer ou non, de manière
pas forcément aléatoire, je ne souhaite surtout pas que les
autres joueurs connaissent mon choix !%
\item Dans une population de joueurs, une strat\'egie
correspond à une proportion dans laquelle une stratégie pure est
jouée dans la population. Par exemple, si dans la ville de New
Delhi 30\% des piétons cherchent à se ranger à gauche
(systématiquement) lorsqu'ils croisent un autre piéton, et 70\% à
droite, à chaque fois que je croise un piéton dans cette ville je
joue face à la stratégie mixte $(30\%,70\%)$.%%
\item Incertitude sur les paiements. (Harsanyi). Supposons que
selon mon ``humeur'' du moment, j'aie une faible préférence pour
bluffer, ou pour ne pas bluffer, et que les autres joueurs ne
connaissent pas mes préférences. Dans ce cas, je joue la stratégie
que je préfère, donc pas aléatoirement, mais les autres joueurs
ont l'impression que je joue aléatoirement, et se représentent mon
comportement par une stratégie mixte.
\end{itemize}

\begin{proposition}
L'ensemble $\Sigma_i$ des stratégies mixtes du joueur $i$ est
convexe. Ses points extr\^emaux sont les stratégies qui mettent
probabilité 1 sur
un seul point de $S_i$. %
\end{proposition}

Remarquons qu'un strat\'egie pure $s_i$ correspond \`a la
strat\'egie mixte qui joue $s_i$ avec probabilit\'e~1. On
considère donc $S_i$ comme sous-ensemble de $\Sigma_i$. Par
conséquent, toute stratégie pure est vue comme une stratégie
mixte. En considérant l'ensemble des stratégies mixtes, on étend
les ensembles des stratégies des joueurs.

\subsection{Extension mixte d'un jeu}

Si chaque joueur $j$ joue la strat\'egie $\sigma_j$, la
probabilité que $(s_j)_j$ soit le profil d'actions effectivement
joué est $\Pi_i \sigma_i(s_i)$. Par conséquence, le paiement
esp\'er\'e du joueur $i$ est :
$$
\sum_{(s_j)\in \prod_j S_j} \Pi_j \sigma_j(s_j) g_i((s_j)_j)
$$
Remarque : On suppose que les fonctions de gain $g_i\colon
S\to\reals$ sont des fonctions d'utilité de von Neumann et
Morgenstern, l'utilité pour l'aléa est l'espérance de l'utilité
obtenue.

La relation précédente définit des fonctions de paiement
$g_i\colon\Sigma\to\reals$. On utilise donc la même notation pour
l'application de $S$ vers $\reals$ et pour son extension de
$\Sigma$ vers $\reals$.

\begin{definition}
{\em L'extension mixte} du jeu sous forme normale $(I,
(S_i),(g_i)_i)$ est le jeu sous forme normale
$(I,(\Sigma_i)_i,(g_i)_i)$.%
\end{definition}

Dans l'extension mixte :
\begin{itemize}%
\item L'ensemble des joueurs est $I$.%
\item Chaque joueur choisit une stratégie mixte $\sigma_i\in
\Sigma_i$. %
\item Le paiement de $i$ est $g_i(\sigma)$.%
\end{itemize}%


\begin{proposition} L'application $g\colon\Sigma\to\reals$ est multilin\'eaire,
c.a.d.\, elle est lin\'eaire sur chaque coordonn\'ee $\Sigma_i$.
\end{proposition}
{\bf Preuve :} il faut montrer que pour
$\sigma_{-j}\in\Sigma_{-j}$, $\sigma_j, \sigma'_j \in \Sigma_j$ et
$\lambda\in[0,1]$, et $\sigma''_j=\lambda \sigma_j + (1-\lambda)
\sigma'_j$:
$$
g(\sigma_{-j},\sigma''_{j}) = \lambda g(\sigma_{-j},\sigma_{j}) +
(1-\lambda)g(\sigma_{-j},\sigma'_{j})
$$
En particulier on a :
$$
g(\sigma_{-j},\sigma_{j}) = \sum_{s_j\in S_j}
\sigma_j(s_j)g(\sigma_{-j},s_j)
$$

\subsection{Équilibre de Nash en stratégies mixtes}

\begin{definition}
 Un \'equilibre de Nash en strat\'egies mixtes
de $G$ est un \'equilibre de Nash du jeu avec ensembles de
strat\'egies $\Sigma_i$ et fonction de paiements $g_i$.
\end{definition}

Un équilibre de Nash en stratégies mixtes est donc un profil de
strat\'egies mixtes $\sigma\in\Sigma$ tel que $\forall i$,
$\forall \sigma'_i\in \Sigma_i$,
$$
g_i(\sigma_{-i},\sigma_{i})\geq g_i(\sigma_{-i},\sigma'_{i})
$$

\section{Etude de l'équilibre de Nash en stratégies mixtes}
Nous allons nous poser (et y répondre) les questions suivantes sur
les EN en stratégies mixtes.\\%
\begin{enumerate}
\item Quel rapport existe-t-il entre les EN en stratégies et en stratégies pures ? %
\item Comment calculer les EN en stratégies mixtes ? %
\item Quel est le rapport entre les EN en stratégies mixtes et l'élimination itérée de stratégies dominées ?%
\item Existence des EN en stratégies mixtes ?
\end{enumerate}%

\subsection{Équilibres de Nash en stratégies pures et mixtes}

\begin{proposition} $\sigma$ est un \'equilibre de Nash en
strat\'egies mixtes si et seulement si :
$$
\forall i,~~\forall s_i\in S_i~~~~~~~
g_i(\sigma_{-i},\sigma_{i})\geq g_i(\sigma_{-i},s_{i})
$$
\end{proposition} La partie ``seulement si'' provient du fait
qu'une strat\'egie pure est aussi une strat\'egie mixte, la partie
``si'' de la lin\'earit\'e de $g_i$ en $\sigma_i$.

\begin{exemple} Chercher un \'equilibre de Nash du jeu de Pile ou Face, puis
du jeu de ``Pierre, Papier, Ciseaux". %
\end{exemple}

\begin{proposition} Tout EN en strat\'egies pures est aussi un EN
en strat\'egies mixtes.
\end{proposition}
{\bf Preuve }: Corollaire de la proposition pr\'ec\'edente.

Par conséquent, considérer les stratégies mixtes nous donne plus
d'équilibres qu'avec les stratégies pures. Pour les équilibres ne
faisant intervenir que des stratégies pures, on parle d'équilibres
de Nash en stratégies pures.

\subsection{Recherche des équilibres en stratégies mixtes}
\begin{definition}
Dans un jeu fini, le support de $\sigma_i\in \Sigma_i$ est
$$
supp(\sigma_i)=\{s_i\in S_i, \sigma_i(s_i)>0\}
$$
\end{definition}

\begin{proposition}(indifférence sur le support)
Si $\sigma$ est un \'equilibre de Nash en strat\'egies mixtes,
$$
\forall i, ~~\forall s_i,s'_i\in
supp(\sigma_i)~~~~~g_i(\sigma_{-i},s_{i})= g_i(\sigma_{-i},s'_{i})
$$
\end{proposition}
En effet, si une stratégie du support donnait un meilleur paiement
que $\sigma_i$, elle constituerait une déviation profitable. Elles
donnent donc toutes un paiement inférieur ou égal à celui de
$\sigma_i$. Mais alors, si une d'entre elles donnait un paiement
strictement inférieur à celui de $\sigma_i$, on obtient une
contradiction (utiliser le fait que cette stratégie est jouée avec
probabilité $> 0$ sous $\sigma_i$ et la linéarité de $g_i$ en
$\sigma_i$).

\begin{proposition} $\sigma$ est un \'equilibre de Nash en
strat\'egies mixtes si et seulement si pour tout $i$, %
\begin{itemize}%
\item $\forall s_i\in supp(\sigma_i)$
$$
g_i(\sigma_{-i},s_{i})= g_i(\sigma_{-i},\sigma_{i})
$$
\item $\forall s_i\not\in supp(\sigma_i)$
$$
g_i(\sigma_{-i},s_{i})\leq g_i(\sigma_{-i},\sigma_{i})
$$
\end{itemize}%
\end{proposition}
Conséquence des résultats précédents.

Ceci fournit une proc\'edure de recherche des \'equilibres :
\begin{enumerate}
\item essayer tous les supports possibles ;%
\item résoudre les probabilités rendant chaque joueur
indifférent sur son
support ;%
\item vérifier que les stratégies hors du support ne donnent pas
un paiement supérieur.
\end{enumerate}

{\bf Exemple} : Quels sont les équilibres de Nash  en stratégies
pures et mixtes du jeu suivant ?
\begin{center}
\begin{game}{2}{2}
     &  $G$    &  $D$ \\
$H$ & $1,1$    & $0,4$\\
$B$ & $0,2$    & $2,1$\\
\end{game}\hspace*{20mm}%
\end{center}
Même question avec le jeu de "Pierre, Papier, Ciseaux".

\subsection{Stratégies mixtes et dominées}

Nous commençons par étendre la relation de domination stricte aux
stratégies mixtes.
\begin{definition}
On dit que  $\sigma_i$ domine $\sigma'_i$ (strictement) lorsque
pour toute $\sigma_{-i}$,
$$
g_i(\sigma_i,\sigma_{-i}) > g_i(\sigma'_i,\sigma_{-i})
$$
\end{definition}

Pour vérifier qu'une stratégie en domine une autre, il suffit de
le vérifier face à tout profil de stratégies pures.
\begin{proposition}
$\sigma_i$ domine $\sigma'_i$ si et seulement si pour tout
$s_{-i}$,
$$
g_i(\sigma_i,s_{-i}) > g_i(\sigma'_i,s_{-i})
$$
\end{proposition}
La condition nécessaire est évidente. Pour la condition
suffisante, utiliser la linéarité.


Une strat\'egie non domin\'ee par une strat\'egie pure peut \^etre
domin\'ee par une strat\'egie mixte.

\begin{exemple}
Dans le jeu suivant :
\begin{center}
\begin{game}{3}{3}
     &  $G$  &  $M$  & $D$\\
$h$ & $1,1$  & $0,2$ & $0,4$\\
$m$ & $0,2$  & $5,0$ & $1,6$\\
$b$ & $0,2$  & $1,1$ & $2,1$\\
\end{game}\hspace*{20mm}%
\end{center}
$M$ est strictement domin\'e par $\frac12G+\frac12D$.%
\end{exemple}

Nous pouvons redéfinir la procédure d'élimination itérée de
stratégies dominées en prenant en compte les stratégies dominées
par des stratégies mixtes.
\begin{itemize}
\item {\bf \'Etape 1 :} $S_i^0 = S_i$.%
\item {\bf \'Etape 2 :} %
$$
S_i^1 = \{\mbox{elts.\,de $S_i$ non dominés par des elts.\,de
$\Sigma_i$}\}
$$%
$$
\Sigma_i^1 = \{\mbox{stratégies mixtes à support dans $S_i^1$}\}
$$
\item {\bf \'Etape $k+1$ :}%
$$%
S_i^{k+1} = \{s_i \in S_i^k, \not\exists\sigma_i \in \Sigma^{k}_i,
\forall s_{-i}\in
S^{k}_{-i}, g_i(\sigma_i,s_{-i})>g_i(s_i,s_{-i})\}$$%
Stratégies non dominées {\em face aux stratégies de } $S^k_i$%
$$
\Sigma_i^k = \{\mbox{stratégies mixtes à support dans $S_i^k$}\}
$$
\item {\bf \'Etape $\infty$ :} $S_i^{\infty} = \cap_k S^k_i$.
\end{itemize}

Les équilibres de Nash en stratégies mixtes résistent à la
procédure d'élimination itérée de stratégies (strictement)
dominées.

\begin{proposition}
Si $\sigma$ est un EN en stratégies mixtes:
$$
\forall i\ \  supp(\sigma_i) \in S_i^\infty
$$
\end{proposition}
Pour la démonstration, suivre le même raisonnement que pour
l'élimination des stratégies dominées par des stratégies pures.
Plus généralement:
\begin{proposition}
Les EN en stratégies mixtes d'un jeu sont les mêmes que celui du
jeu duquel on a éliminé les stratégies strictement
dominées.
\end{proposition}

{\bf Exemple} : r\'esoudre le jeu pr\'ec\'edent.

\subsection{Existence des équilibres de Nash en stratégies mixtes}
On a un équilibre de Nash lorsque chaque joueur joue optimalement
face aux stratégies des autres joueurs.  Pour caractériser ces
équilibres, nous sommes intéressés par l'application qui associe à
chaque $\sigma_{-i}$ l'ensemble des stratégies $\sigma_i$ qui
donnent le meilleur paiement face à $\sigma_{-i}$. Cette
correspondance s'appelle la correspondance de meilleures réponses.

\begin{definition}
Une correspondance de $A$ vers $B$ est une application de $A$ vers
les parties de $B$.%
\end{definition}

\begin{definition}
 La correspondance de meilleure r\'eponse du joueur $i$ est la correspondance de
$\Sigma_{-i}$ vers $\Sigma_i$ donnée par :
$$
MR_i(\sigma_{-i})=\arg\max_{\sigma_i\in
\Sigma_i}g_i(\sigma_{-i},\sigma_i)
$$
La correspondance de meilleures réponses du jeu $G$ est la
correspondance de $\Sigma$ vers $\Sigma$ donnée par
$$
MR(\sigma) = \Pi_i MR_i(\sigma_{-i})
$$
\end{definition}
Les meilleures réponses permettent de caractériser les équilibres
de Nash :

\begin{proposition}
$\sigma$ est un EN si et seulement si :
$$
\forall i\ \  \sigma_i \in MR_i(\sigma_{-i})
$$
i.e.
$$
\sigma \in MR(\sigma)
$$
\end{proposition}

On a les propriétés suivantes de la correspondance de meilleures
réponses :

\begin{proposition}
\begin{itemize}
\item Pour tout $\sigma$, $MR(\sigma)$ est non vide%
\item Pour tout $\sigma$, $MR(\sigma)$ est convexe%
\item Si tous les $S_i$ sont finis, le graphe de $MR$ est fermé,
i.e.\,si $\sigma_n\to \sigma$, $\tau_n\to \tau$ et $\tau_n\in
MR(\sigma_n)$ alors $\tau\in MR(\sigma)$.%
\end{itemize}
\end{proposition}

On rappelle le théorème d'existence de point fixe d'une
correspondance de Kakutani :
\begin{theoreme}[de Kakutani] Soit $X$
un compact convexe non vide de $\reals^n$, et $F$ une
correspondance de $X$ vers $X$ telle que :
\begin{itemize}
\item $\forall x\in X$, $F(x)$ est convexe et non vide ; \item Le
graphe de $F$ est ferm\'e.
\end{itemize}
Alors $F$ admet un point fixe, i.e.\, il existe $x\in X$ tel que
$x\in F(x)$. %
\end{theoreme}

{\bf Exercice} Montrer que les hypothèses
\begin{itemize}
\item $X$ compact ;%
\item $X$ convexe ; \item $F(x)$ non vide pour
tout $x$ ;%
\item $F(x)$ convexe pour tout $x$ ;%
\item le graphe de $F$ est ferm\'e
\end{itemize}
sont nécessaires.

En corollaire du théorème précédent, on a le théorème d'existence
d'équilibres de Nash en stratégies mixtes :


\begin{theoreme}[Glicksberg, Nash] Tout jeu fini admet un
équilibre de Nash en stratégies mixtes.
\end{theoreme}

\section{Le jeu de p\'enalties}

Un joueur d\'ecide de tirer soit \`a droite, soit \`a gauche du
gardien. Le gardien peut plonger soit \`a droite, soit \`a gauche.
Le tir peut \^etre bien tir\'e ou sortir complètement (avec
probabilit\'e $p_D$ pour un tir \`a droite, $p_G$ \`a gauche). Le
gardien arr\^ete toujours la balle s'il plonge du bon c\^ot\'e. Il
y a donc but si :
\begin{enumerate}
\item La balle est bien tir\'ee ; %
\item Le gardien plonge du mauvais c\^ot\'e.
\end{enumerate}

On demande de résoudre les questions suivantes : L'objectif du
joueur est de maximiser les possibilit\'es de marquer but, et le
gardien veut minimiser ces possibilit\'es.
\begin{itemize}
\item Repr\'esenter le jeu sous forme normale. %
\item Quels sont les \'equilibres de Nash ? %
\item Un droitier tire-t-il plus souvent \`a gauche ou \`a droite
? %
\item Face \`a un droitier, le gardien plonge-t-il plus souvent
\`a gauche ou \`a droite ?
\end{itemize}
Indice : Pour un droitier, $p_G > p_D$.

\section{Course au brevet} Deux entreprises se lancent dans une course au
brevet. Chacune investit un montant positif dans la recherche et
celui qui investit le plus reçoit une subvention $V>0$. En cas
d'\'egalit\'e, chacun reçoit $V/2$. On prend l'intervalle $\lbrack
0,V\rbrack$ comme ensemble de strat\'egies pour chaque joueur.

\begin{enumerate}
\item Montrer que ce jeu n'a pas d'\'equilibre (en strat\'egies
pures). %
\item Montrer qu'il existe un \'equilibre en strat\'egies mixtes
dans lequel joue une loi \`a support ``plein'' $\lbrack
0,V\rbrack$.\
\end{enumerate}

\chapter{Jeux sous forme extensive}

Les jeux sous forme normale donne une représentation adéquate de
joueurs effectuant des choix de stratégies simultanés. Dans ces
jeux, aucun jouer n'a d'information sur les stratégies utilisées
par les autres au moment de faire ses choix. Dans de nombreux cas
concrets, les joueurs effectuent leurs choix en fonction d'une
situation observée qui dépend en particulier des choix précédents
des autres joueurs. Par exemple c'est le cas pour une enchère
croissante, dans une négociation, ou dans le jeu d'échecs...

Pour représenter ces possibilités, nous allons généraliser la
représentation sous forme d'arbres et les  ensembles d'information
vus au chapitre de la décision dynamique. La seule différence est
que nous avons plusieurs joueurs là où nous n'avions qu'un seul
agent.

\section{Exemples préliminaires}
Nous traitons une série d'exemples avant de définir la notion de
jeu sous forme extensive.

\subsection{Jeu d'entrée} Une firme $M$ est en situation de
monopole sur un marché. Une autre firme $E$ peut décider d'entrer
($e$) ou non ($n$) sur le marché. Si la firme $E$ décide d'entrer,
la firme $M$ peut alors soit combattre ($c$), soit s'accommoder
($a$). Les paiements pour $M$ et pour $A$ sont $(2,0)$ si $E$
n'entre pas, $(-1,-1)$ si $E$ entre et $M$ combat, et $(1,1)$ si
$E$ entre et $M$ s'accommode. On représente ce jeu par l'arbre suivant:\\[1cm]
\begin{tabular}{c@{\hspace{1cm}}c@{\hspace{1cm}}c@{\hspace{1cm}}c@{\hspace{1cm}}}
&&\circlenode{a}{\makebox[0.2cm][c]{\small $E$}}&\\[0.7cm]
&\circlenode{b}{\makebox[0.2cm][c]{\small $M$}}&&\rnode{c}{\small$2$}\\
&&&{\small$0$}\\[0.4cm]
\rnode{d}{\small$-1$}&&\rnode{e}{\small$1$}&\\
{\small$-1$}&&{\small$1$}&\\[0.7cm]
\ncline[linewidth=1pt,linecolor=yellow]{-}{a}{b}\bput[1pt]{0}(0.5){\small$e$}%
\ncline[linewidth=1pt,linecolor=yellow]{-}{a}{c}\aput[1pt]{0}(0.5){\small$n$}%
\ncline[linewidth=1pt,linecolor=yellow]{-}{b}{d}\bput[1pt]{0}(0.5){\small$c$}%
\ncline[linewidth=1pt,linecolor=yellow]{-}{b}{e}\aput[1pt]{0}(0.5){\small$a$}%
%\rput{45}{\psellipse[linecolor=green](4.2,1)(1.3,0.3)}
\end{tabular}

\subsection{Compétition de Stackelberg}
Supposons qu'une firme, 1, ait la possibilité de choisir son
niveau de production $q_1$ avant que la firme $2$ ne fixe le sien
$q_2$. Pour fixer les idées, supposons que $q_i \in \{0,1,2\}$,
que la fonction de demande est $p=\max\{3-q_1+q_2,0\}$, et que le
co\^ut de production est nul pour chaque firme.
On représente ce jeu par l'arbre suivant:\\[0.8cm]
\begin{center}
\begin{tabular}{c@{\hspace{1cm}}c@{\hspace{1cm}}c@{\hspace{1cm}}c@{\hspace{1cm}}c@{\hspace{1cm}}c@{\hspace{1cm}}%
c@{\hspace{1cm}}c@{\hspace{1cm}}c@{\hspace{1cm}}}
&&&&\circlenode{a}{\makebox[0.2cm][c]{\small $1$}}&&&&\\[1cm]
&\circlenode{b}{\makebox[0.2cm][c]{\small $2$}}&&&
\circlenode{c}{\makebox[0.2cm][c]{\small $2$}}&&&
\circlenode{d}{\makebox[0.2cm][c]{\small $2$}}&\\[1cm]
\rnode{e}{\small$0$}&\rnode{f}{\small$0$}&\rnode{g}{\small$0$}&%
\rnode{h}{\small$2$}&\rnode{i}{\small$1$}&\rnode{j}{\small$0$}&%
\rnode{k}{\small$2$}&\rnode{l}{\small$0$}&\rnode{m}{\small$0$}\\
\small$0$&\small$2$&\small$2$&%
\small$0$&\small$1$&\small$0$&%
\small$0$&\small$0$&\small$0$\\%
\ncline[linewidth=1pt,linecolor=yellow]{-}{a}{b}\bput[1pt]{0}(0.5){\small$0$}%
\ncline[linewidth=1pt,linecolor=yellow]{-}{a}{c}\aput[1pt]{0}(0.5){\small$1$}%
\ncline[linewidth=1pt,linecolor=yellow]{-}{a}{d}\aput[1pt]{0}(0.5){\small$2$}%
\ncline[linewidth=1pt,linecolor=yellow]{-}{b}{e}\bput[1pt]{0}(0.5){\small$0$}%
\ncline[linewidth=1pt,linecolor=yellow]{-}{b}{f}\bput[1pt]{0}(0.5){\small$1$}%
\ncline[linewidth=1pt,linecolor=yellow]{-}{b}{g}\aput[1pt]{0}(0.5){\small$2$}%
\ncline[linewidth=1pt,linecolor=yellow]{-}{c}{h}\bput[1pt]{0}(0.5){\small$0$}%
\ncline[linewidth=1pt,linecolor=yellow]{-}{c}{i}\bput[1pt]{0}(0.5){\small$1$}%
\ncline[linewidth=1pt,linecolor=yellow]{-}{c}{j}\aput[1pt]{0}(0.5){\small$2$}%
\ncline[linewidth=1pt,linecolor=yellow]{-}{d}{k}\bput[1pt]{0}(0.5){\small$0$}%
\ncline[linewidth=1pt,linecolor=yellow]{-}{d}{l}\bput[1pt]{0}(0.5){\small$1$}%
\ncline[linewidth=1pt,linecolor=yellow]{-}{d}{m}\aput[1pt]{0}(0.5){\small$2$}%
\end{tabular}
\end{center}

\subsection{Pile ou face}

Le joueur 1 choisit pile ou face, puis le joueur 2 choisit pile ou
face. Le joueur 2 gagne si les choix sont égaux, et perd si les
choix sont différents (c'est un jeu à somme nulle).
On représente ce jeu par l'arbre suivant:\\[0.5cm]
\begin{center}
\begin{tabular}{c@{\hspace{0.5cm}}c@{\hspace{0.5cm}}c@{\hspace{0.5cm}}%
c@{\hspace{0.5cm}}c@{\hspace{0.5cm}}c@{\hspace{0.5cm}}%
c@{\hspace{0.5cm}}}
&&&\circlenode{a}{\makebox[0.2cm][c]{\small $1$}}&&&\\[1cm]
&\circlenode{b}{\makebox[0.2cm][c]{\small $2$}}&&&&
\circlenode{c}{\makebox[0.2cm][c]{\small $2$}}&\\[1cm]
\rnode{d}{\small$1$}&&\rnode{e}{\small$-1$}&&\rnode{f}{\small$-1$}%
&&\rnode{g}{\small$1$}\\
\small$-1$&&\small$1$&&\small$1$&&\small$-1$\\%
\ncline[linewidth=1pt,linecolor=yellow]{-}{a}{b}\bput[1pt]{0}(0.5){\small$P$}%
\ncline[linewidth=1pt,linecolor=yellow]{-}{a}{c}\aput[1pt]{0}(0.5){\small$F$}%
\ncline[linewidth=1pt,linecolor=yellow]{-}{b}{d}\bput[1pt]{0}(0.5){\small$P$}%
\ncline[linewidth=1pt,linecolor=yellow]{-}{b}{e}\aput[1pt]{0}(0.5){\small$F$}%
\ncline[linewidth=1pt,linecolor=yellow]{-}{c}{f}\bput[1pt]{0}(0.5){\small$P$}%
\ncline[linewidth=1pt,linecolor=yellow]{-}{c}{g}\aput[1pt]{0}(0.5){\small$F$}%
\end{tabular}
\end{center}

\subsection{Pile ou Face sans observation}
On suppose maintenant que 2 ne connaît pas le choix de 1 lorsqu'il
effectue son choix. Ceci peut être représenté à l'aide d'un
ensemble d'information :
\begin{center}
\begin{tabular}{c@{\hspace{0.5cm}}c@{\hspace{0.5cm}}c@{\hspace{0.5cm}}%
c@{\hspace{0.5cm}}c@{\hspace{0.5cm}}c@{\hspace{0.5cm}}%
c@{\hspace{0.5cm}}}
&&&\circlenode{a}{\makebox[0.2cm][c]{\small $1$}}&&&\\[1cm]
&\circlenode{b}{\makebox[0.2cm][c]{\small $2$}}&&&&
\circlenode{c}{\makebox[0.2cm][c]{\small $2$}}&\\[1cm]
\rnode{d}{\small$1$}&&\rnode{e}{\small$-1$}&&\rnode{f}{\small$-1$}%
&&\rnode{g}{\small$1$}\\
\small$-1$&&\small$1$&&\small$1$&&\small$-1$\\%
\ncline[linewidth=1pt,linecolor=yellow]{-}{a}{b}\bput[1pt]{0}(0.5){\small$P$}%
\ncline[linewidth=1pt,linecolor=yellow]{-}{a}{c}\aput[1pt]{0}(0.5){\small$F$}%
\ncline[linewidth=1pt,linecolor=yellow]{-}{b}{d}\bput[1pt]{0}(0.5){\small$P$}%
\ncline[linewidth=1pt,linecolor=yellow]{-}{b}{e}\aput[1pt]{0}(0.5){\small$F$}%
\ncline[linewidth=1pt,linecolor=yellow]{-}{c}{f}\bput[1pt]{0}(0.5){\small$P$}%
\ncline[linewidth=1pt,linecolor=yellow]{-}{c}{g}\aput[1pt]{0}(0.5){\small$F$}%
\ncline[linewidth=2pt,linecolor=green,linestyle=dotted]{-}{b}{c}% ensemble d'information
\end{tabular}
\end{center}
L'ensemble d'information du joueur 2 comprend ses deux n{\oe}uds
de décision. Ceci signifie que le joueur 2 ne différencie pas ces
deux n{\oe}uds. Il doit effectuer un choix indépendant de celui
effectué précédemment par le joueur 1.
\subsection{Pile ou Face contre la nature} Cette fois ci, il
n'y a pas de joueur 2, mais le joueur 1 joue contre la Nature qui
tire une pièce de monnaie, Pile ou Face, avec probabilités $\lambda$ et $1-\lambda$.\\[0.5cm]
\begin{center}
\begin{tabular}{c@{\hspace{0.5cm}}c@{\hspace{0.5cm}}c@{\hspace{0.5cm}}%
c@{\hspace{0.5cm}}c@{\hspace{0.5cm}}c@{\hspace{0.5cm}}%
c@{\hspace{0.5cm}}}
&&&\circlenode{a}{\makebox[0.2cm][c]{\small $\red N$}}&&&\\[1cm]
&\circlenode{b}{\makebox[0.2cm][c]{\small $1$}}&&&&
\circlenode{c}{\makebox[0.2cm][c]{\small $1$}}&\\[1cm]
\rnode{d}{\small$1$}&&\rnode{e}{\small$0$}&&\rnode{f}{\small$0$}%
&&\rnode{g}{\small$1$}\\
\ncline[linewidth=1pt,linecolor=yellow]{-}{a}{b}\bput[1pt]{0}(0.6){\small$P$}%
\ncline[linewidth=1pt,linecolor=yellow]{-}{a}{c}\aput[1pt]{0}(0.6){\small$F$}%
\ncline[linewidth=1pt,linecolor=yellow]{-}{a}{b}\bput[1pt]{0}(0.3){\small\red$\lambda$}%
\ncline[linewidth=1pt,linecolor=yellow]{-}{a}{c}\aput[1pt]{0}(0.3){\small\red$1-\lambda$}%
\ncline[linewidth=1pt,linecolor=yellow]{-}{b}{d}\bput[1pt]{0}(0.5){\small$P$}%
\ncline[linewidth=1pt,linecolor=yellow]{-}{b}{e}\aput[1pt]{0}(0.5){\small$F$}%
\ncline[linewidth=1pt,linecolor=yellow]{-}{c}{f}\bput[1pt]{0}(0.5){\small$P$}%
\ncline[linewidth=1pt,linecolor=yellow]{-}{c}{g}\aput[1pt]{0}(0.5){\small$F$}%
\ncline[linewidth=2pt,linecolor=green,linestyle=dotted]{-}{b}{c}% ensemble d'information
\end{tabular}
\end{center}
$N$ représente la nature. On a représenté les probabilités
$\lambda, 1-\lambda$ des choix pouvant être faits par la nature.

\section{Jeu sous forme extensive $\Gamma$}

Voici la définition formelle d'un jeu sous forme extensive. On ne
demandera pas de la connaître par c{\oe}ur. Cependant, il est
important de savoir reconnaître un tel jeu et de savoir
représenter une situation stratégique par un jeu sous forme
extensive.
\begin{definition}
Un jeu sous forme extensive $\Gamma$ est donné par :

\begin{itemize}
\item $I$ ensemble fini de joueurs,  $N\not\in I$ (nature);
\item $T$ ensemble fini de n{\oe}uds, $Z\subset T$ n{\oe}uds terminaux, et pour
$t\not\in Z$ :
\begin{itemize}
\item Le joueur ou la nature $i(t)\in \{N\}\cup I$ qui joue en $t$ ;
\item L'ensemble de choix possibles $A(t)$ ;
\item Le n{\oe}ud successeur résultant du choix $a$, $N(t,a)$ ;
\item Une distribution $D(t)$ sur $A(t)$ si $i(t)= N$.
\end{itemize}
\item $u_i:Z\to\reals$ fonctions d'utilité de $i$ ;
\item $h(t)$ classe d'équivalence de $t$ pour l'information :
$$
x'\in h(x) \implies i(x')= i(x),~A(x')=A(x), h(x')=h(x)
$$
\end{itemize}

Afin que le jeu ait une structure d'arbre, on requiert que les
éléments vérifient les propriétés suivantes.

\begin{itemize}
\item $s(t)=\{N(t,a),a\in A(t)\}$ ensemble des successeurs directs de $t$. ($\emptyset$ si $t\in Z$).
\item $S(t)= s(t)\cup s(s(t)) \cup s(s(s(t)))...$ tous les successeurs de $t$.
\end{itemize}
\begin{itemize}
\item Il n'y a pas de cycle : $\forall t, t\not\in S(t)$
\item N{\oe}ud initial $\exists\tilde t, S(\tilde t)\cup\{\tilde t\}=T$
\item Chaque n{\oe}ud non initial a un unique prédécesseur:
$\forall t\neq \tilde t,~\exists ! t',~t\in s(t')$
\end{itemize}
\end{definition}
En exercice, on pourra définir les composants $T$, $Z$, etc qui
représentent les jeux précédents, puis vérifier que ces jeux ont
bien les propriétés de structure d'arbre.

\subsection{Hypothèse de mémoire parfaite}

Il est naturel de supposer qu'un joueur ne peut oublier de
l'information au cours du déroulement du jeu. Notamment, ce joueur
doit se rappeler tous ses choix précédents. C'est l'hypothèse dite
de mémoire parfaite.

\begin{definition}
Le jeu $\Gamma$ est à mémoire parfaite si pour tout $x, y, z$ tels
que $i(x)=i(y)=i(z)$, $h(y)=h(z)$ et $a$ tels que $z\in
S(N(x,a))\cup\{N(x,a)\}$, il existe $w$ tel que :
\begin{itemize}
\item $h(w)=h(x)$ ;
\item $y\in S(N(w,a))\cup\{N(w,a)\}$.
\end{itemize}
\end{definition}

A titre d'exemple, considérons le cas du conducteur distrait. Une
autoroute comprend deux sorties. Afin de se rendre chez lui, un
automobiliste doit prendre la seconde sortie. Cependant, on
suppose que le conducteur est distrait, c'est-à-dire qu'au moment
de décider de sortir de l'autoroute ou d'y rester, il est
incapable de déterminer s'il se trouve à la première ou a la
deuxième sortie. Sauriez-vous représenter cette situation par un
arbre de jeu ? L'hypothèse de mémoire parfaite est-elle vérifiée ?
Pourquoi ?

Sauf exception, on ne considérera des jeux avec mémoire parfaite. %

\subsection{Stratégies pures}
Comme dans le cas de la décision dynamique, une stratégie pour un
joueur est une fonction de ces ensembles d'information vers ses
choix.

\begin{definition}
 Une stratégie pure pour $i\in I$ dans $\Gamma$ est une
application $s_i\colon \{t,i(t)=i\}\to\cup_{t}A(t)$ telle que:
\begin{itemize}
\item $\forall t, s_i(t)\in A(t)$ ;
\item $s_i(t) = s_i(t')$ si $t\in h(t')$.
\end{itemize}
\end{definition}

Remarquons que cette définition permet facilement d'énumérer les
stratégies du joueur. Cependant, deux stratégies peuvent toujours
donner le même résultat. On dira alors qu'elles sont équivalentes
(nous reviendrons plus tard sur cette notion).

Comme exercice, et à titre de rappel des problèmes de décision
dynamiques, quelles sont les stratégies du joueur 1 dans le jeu
suivant ? Quelles stratégies sont équivalentes ?

\begin{tabular}{c@{\hspace{1cm}}c@{\hspace{1cm}}c@{\hspace{1cm}}c@{\hspace{1cm}}}
&&\circlenode{a}{\makebox[0.2cm][c]{\small $1$}}&\\[0.7cm]
&\circlenode{b}{\makebox[0.2cm][c]{\small $1$}}&&\rnode{c}{\small$2$}\\
&&&{}\\[0.4cm]
\rnode{d}{\small$-1$}&&\rnode{e}{\small$1$}&\\
{}&&{}&\\[0.7cm]
\ncline[linewidth=1pt,linecolor=yellow]{-}{a}{b}\bput[1pt]{0}(0.5){\small$a$}%
\ncline[linewidth=1pt,linecolor=yellow]{-}{a}{c}\aput[1pt]{0}(0.5){\small$b$}%
\ncline[linewidth=1pt,linecolor=yellow]{-}{b}{d}\bput[1pt]{0}(0.5){\small$c$}%
\ncline[linewidth=1pt,linecolor=yellow]{-}{b}{e}\aput[1pt]{0}(0.5){\small$d$}%
%\rput{45}{\psellipse[linecolor=green](4.2,1)(1.3,0.3)}
\end{tabular}

\subsection{Stratégies aléatoires}

Comme dans le cas des jeux sous forme normale, nous avons besoin
de représenter les possibilités de jouer de manière aléatoire.
Deux types de stratégies représentent des choix aléatoires des
joueurs.
\begin{definition}
Une stratégie mixte est une loi de probabilité sur les stratégies
pures.
\end{definition}


\begin{definition}
Une stratégie de comportement est une application $\sigma_i\colon
\{t,i(t)=i\}\to\cup_{t}\Delta(A(t))$ telle que :
\begin{itemize}
\item $\forall t, \sigma_i(t)(a)>0\implies a\in A(t)$ ;
\item $\sigma_i(t) = \sigma_i(t')$ si $t\in h(t')$.
\end{itemize}
\end{definition}

Dans une stratégie mixte, le joueur choisit aléatoirement ``au
début du jeu'' la stratégie pure qu'il utilisera par la suite. Une
fois qu'il a choisi cette stratégie, il joue en suivant cette
règle déterministe de décisions. Une stratégie de comportements,
spécifie un choix aléatoire à chaque n{\oe}ud de décision, on donc
un aléa a chaque fois qu'il s'agit de prendre une décision. L'aléa
porte sur la prochaine action choisie, et non pas sur une règle
globale de comportement.

\subsection{\'Equivalence des stratégies mixtes et de comportement}

Deux stratégies $\sigma_i$ et $\sigma_i'$ du joueur $i$ (pures,
mixtes ou de comportement) sont équivalentes lorsque pour toutes
stratégies $\sigma_{-i}$ des autres joueurs (idem),
$(\sigma_i,\sigma_{-i})$ et $(\sigma'_i,\sigma_{-i})$ induisent la
m\^eme distribution sur les n{\oe}uds terminaux.

Comme nous l'avons vu, des stratégies pures peuvent en particulier
être équivalentes entre elles.

\subsection{Théorème de Kuhn}

\begin{theoreme}[de Kuhn]
Dans tout jeu à mémoire parfaite,
\begin{enumerate}
\item Toute stratégie mixte admet une stratégie de comportement
qui lui est équivalente ;
\item Toute stratégie de comportement
admet une stratégie mixte qui lui est équivalente.
\end{enumerate}
\end{theoreme}

Les deux manières de modéliser la capacité de jouer aléatoirement
-- en utilisant les stratégies mixtes ou les stratégies de
comportement -- sont donc équivalentes dans les jeux à mémoire
parfaite. On utilisera en pratique la représentation la plus
commode selon les cas.%

Ce résultat n'est vrai que pour les jeux à mémoire parfaite.
Reprenez l'exemple du conducteur distrait. Toutes les stratégies
mixtes sont-elles équivalentes à des stratégies de comportement ?
Le contraire est-il vrai ?

\section{Forme normale et équilibres de Nash}
Pour ``résoudre'' des jeux en forme extensive, nous allons
utiliser le concept d'équilibre de Nash de la forme normale, puis
montrer que la forme extensive permet de faire des prédictions
plus fines (équilibre parfait dans les sous-jeux).

\subsection{Forme normale et forme normale réduite}
La forme normale associée à
$\Gamma$ est le jeu sous forme normale $G = (I,(S_i)_i,(g_i)_i)$,
avec
$$
g_i(s) = \E_{P_s} u_i(z)
$$
où $P_s$ est la probabilité sur les n{\oe}uds terminaux induite
par le profil de stratégies (pures) $s$.

Soit $S^r_i$ l'ensemble des classes d'équivalences de stratégies
du joueur $i$. Un élément de $S^r_i$ est appelé
une stratégie réduite du joueur $i$.\\
La forme normale réduite de $\Gamma$ est le jeu sous forme normale
$G^r = (I,(S^r_i)_i,(g^r_i))$ avec
$$
g^r_i (s^r)= g_i(s)
$$
où $s^r = (s^r_j)_j$, $s = (s_j)_j$, $s_j\in s^r_j$ pour tout $j$.
(Le choix de $s_j$ dans $s_j^r$ est indifférent.)

\subsection{\'Equilibre de Nash}

\begin{definition}
Un équilibre de Nash d'un jeu sous forme extensive est un
équilibre de Nash de la forme normale associée.
\end{definition}

{\bf Remarque :} Tout les équilibres de Nash d'un jeu sont obtenus
dans la forme normale réduite.

{\bf Remarque :} Lorsqu'on recherche les équilibres de Nash d'un
jeu sous forme extensive, on utilise la notion de stratégies
mixtes. Cette notion est ici plus utile que celle de stratégies de
comportement.

{\bf Exemple :}
Quels sont les équilibres de Nash du jeu d'entrée (en stratégies pures ou mixtes)~?\\
\begin{center}
\begin{tabular}{c@{\hspace{1cm}}c@{\hspace{1cm}}c@{\hspace{1cm}}c@{\hspace{1cm}}}
&&\circlenode{a}{\makebox[0.2cm][c]{\small $E$}}&\\[0.7cm]
&\circlenode{b}{\makebox[0.2cm][c]{\small $M$}}&&\rnode{c}{\small$2$}\\
&&&{\small$0$}\\[0.4cm]
\rnode{d}{\small$-1$}&&\rnode{e}{\small$1$}&\\
{\small$-1$}&&{\small$1$}&\\[0.7cm]
\ncline[linewidth=1pt,linecolor=yellow]{-}{a}{b}\bput[1pt]{0}(0.5){\small$E$}%
\ncline[linewidth=1pt,linecolor=yellow]{-}{a}{c}\aput[1pt]{0}(0.5){\small$N$}%
\ncline[linewidth=1pt,linecolor=yellow]{-}{b}{d}\bput[1pt]{0}(0.5){\small$C$}%
\ncline[linewidth=1pt,linecolor=yellow]{-}{b}{e}\aput[1pt]{0}(0.5){\small$A$}%
%\rput{45}{\psellipse[linecolor=green](4.2,1)(1.3,0.3)}
\end{tabular}
\end{center}

\section{Jeux à information parfaite et complète}

\begin{definition}
C'est un jeu sous forme extensive dans lequel chaque ensemble
d'information est un singleton.
\end{definition}
Des exemples sont fournis par jeu de Nim, d'échecs, de Go, et par
le jeu d'entrée précédent.

Revenons sur ce jeu d'entrée. La menace de jouer $c$ si $E$ joue $e$ est-elle crédible~?\\
\begin{itemize}
\item Supposons que $E$ ait déjà joué $e$. $M$ a-t-il intérêt
à jouer $c$ ou $a$~? ($E$ est ``mis devant le fait accompli'').\\
\item Ayant effectué ce raisonnement, $E$ a-t-il intérêt à jouer
$e$ ou bien $s$~?
\end{itemize}

\subsection{Procédure d'induction à rebours} Comme pour la
programmation dynamique (jeux à un seul joueur). On définit les
stratégies des joueurs en partant des n{\oe}uds terminaux, puis
pour les décisions plus ``en amont'' de l'arbre progressivement
jusqu'au n{\oe}ud initial.

On obtient une stratégie unique pour chaque joueur si pas
d'égalité entre deux paiements de n{\oe}uds terminaux et s'il n'y
a pas de n{\oe}uds où joue la nature.

\begin{theoreme}[dit de Zermelo]
Tout jeu à information complète et parfaite admet un équilibre de
Nash en stratégies pures, qui peut \^etre calculé par la procédure
d'induction à rebours.
\end{theoreme}

 Application au jeu d'échecs : Dans le cas du jeu d'échecs,
appliquer le théorème de Zermelo donne le résultat suivant :
\begin{itemize}
\item Soit un joueur a une stratégie gagnante ;
\item Soit chacun des joueurs peut se garantir la nulle.
\end{itemize}
Pourquoi joue-t-on encore aux échecs ?

\section{\'Equilibre parfait dans les sous-jeux}
La procédure d'induction à rebours dans les jeux à information
complète et parfaite permet de résoudre les jeux en ayant des
prédictions plus fortes que celles de l'équilibre de Nash. Cette
procédure et ces raisonnements peuvent-ils être généralisés aux
jeux où les ensembles d'information ne sont pas des singletons ?

Nous commençons par définir la notion de sous-jeu, qui peut-être
vu comme une partie d'un jeu pouvant être étudiée séparément du
jeu dans son entièreté.

\begin{definition} Un sous-jeu $\Gamma'$ de $\Gamma$ consiste en
\begin{itemize}
\item Un sous-ensemble de n{\oe}uds non terminaux $T'\subset T$ tel que $T' = S(t')$ pour
un certain $t'$ et $\forall t\in T'$, $h(t)\subset T'$.
\item Les joueurs, mouvements de la nature, actions possibles, successeurs,
utilités et ensembles d'information de $\Gamma'$ sont définis
comme dans $\Gamma$.
\end{itemize}
\end{definition}

{\bf Remarque} : En particulier, $\Gamma$ est un sous-jeu de
$\Gamma$. Tout jeu admet par conséquent au moins un sous-jeu :
lui-même. Les sous-jeux de $\Gamma$ non égaux à $\Gamma$ sont
appelés sous-jeux stricts de $\Gamma$.

{\bf Remarque} : Dans un jeu à information parfaite et complète,
on a autant de sous-jeux que de n{\oe}uds non terminaux..

Si on a des stratégies de comportement pour les joueurs dans un
jeu, on connaît leurs probabilités de choix à tous les n{\oe}uds.
Ceci définit en particulier des stratégies dans tous les
sous-jeux.

\begin{definition} Un profil de stratégies $s$ (pures ou de
comportement) est un équilibre parfait dans les sous-jeux de
$\Gamma$ si $s$ induit un équilibre de Nash dans tout les
sous-jeux de $\Gamma$.
\end{definition}

La proposition suivante est immédiate, car le jeu tout entier est
toujours un sous-jeu de lui-même.

\begin{proposition} Un équilibre parfait dans les sous-jeux est
un équilibre de Nash de $\Gamma$. %
\end{proposition}

\begin{proposition}
Tout jeu $\Gamma$ admet un équilibre de Nash parfait dans les
sous-jeux.
\end{proposition}

{\bf Méthode} : Appliquer la procédure d'induction à rebours à
chaque fois que l'on peut.
\begin{enumerate}
\item Identifier un sous-jeu ``minimal'' $\Gamma'$;
\item Calculer un équilibre de Nash de $\Gamma'$ ;
\item Remplacer $\Gamma'$ par ce paiement de Nash ;
\item Recommencer en $1$ tant que possible.
\end{enumerate}

\end{document}
